% Układ roboczy do redakcji treści. Nie jest zatwierdzonym szablonem PW.
\documentclass[12pt,a4paper,oneside]{report}

\usepackage{fontspec}
\setmainfont{Latin Modern Roman}
\setsansfont{Latin Modern Sans}
\setmonofont{Latin Modern Mono}
\usepackage{polyglossia}
\setdefaultlanguage{polish}
\usepackage[margin=28mm]{geometry}
\usepackage{microtype}
\usepackage{amsmath,amssymb}
\usepackage{graphicx}
\usepackage{booktabs,tabularx}
\usepackage[backend=biber,style=numeric,sorting=nyt,giveninits=true,maxbibnames=99]{biblatex}
\addbibresource{bibliography/references.bib}
\usepackage[hidelinks]{hyperref}

\newcolumntype{Y}{>{\raggedright\arraybackslash}X}
\newcolumntype{P}[1]{>{\raggedright\arraybackslash}p{#1}}
\renewcommand{\arraystretch}{1.2}
\setlength{\emergencystretch}{2em}
% Podmiana ramki na zewnętrzny plik nie wymaga zmiany podpisu ani etykiety.
\newcommand{\plannedfigure}[4]{%
    \begin{figure}[tbp]
        \centering
        \IfFileExists{#1}{%
            \includegraphics[width=\linewidth,height=0.65\textheight,keepaspectratio]{#1}%
        }{%
            \fbox{%
                \begin{minipage}{0.90\linewidth}
                    \centering
                    \vspace{0.8em}{\small #4}\par\vspace{0.8em}
                \end{minipage}%
            }%
        }
        \caption{#2}
        \label{#3}
    \end{figure}%
}


\title{Głębokie sieci neuronowe w zastosowaniu do klasyfikacji danych tabelarycznych\\[1em]\large Wersja robocza rozdziałów 1--5 i~szkielet części wynikowej}
\author{}
\date{6 listopada 2026}

\begin{document}
\maketitle
\tableofcontents

\clearpage
\section{Wstęp}
\label{chap:wstep}

\subsection{Motywacja i znaczenie problemu}
\label{sec:motywacja}

Dane opisujące transakcje, pacjentów, klientów lub wyniki pomiarów często mają postać tabeli, w~której każdej obserwacji odpowiada wiersz, a~poszczególnym właściwościom --- kolumny. W~zadaniu klasyfikacji na podstawie takiego opisu należy przewidzieć etykietę należącą do skończonego zbioru klas. Przykładami są wykrywanie nadużyć, ocena ryzyka kredytowego i~przewidywanie ponownej hospitalizacji. Wspólna organizacja danych nie oznacza jednak jednorodności problemów: kolumny mogą mieć różne typy, skale i~znaczenie, a~obserwacje mogą być ze sobą powiązane~\cite{borisov2022survey}.

Zastosowanie sieci neuronowej do danych tabelarycznych wymaga uzasadnienia zarówno jakością predykcji, jak i~kosztem obliczeń. Punktem odniesienia są metody liniowe, klasyfikatory sąsiedztwa oraz zespoły drzew. Na wynik porównania wpływają także przygotowanie danych, ustawienia uczenia, dobór hiperparametrów i~procedura walidacji. Dlatego ocena architektury wymaga odpowiednio dobranych modeli bazowych oraz spójnych warunków eksperymentu~\cite{gorishniy2021revisiting}.

W~pracy skupiono się na trzech trudnościach: małej liczbie przykładów uczących, dużej liczbie cech w~stosunku do liczby przykładów oraz niezrównoważeniu klas. W~ostatnim przypadku wysoki odsetek poprawnych predykcji może współwystępować z~niską wykrywalnością zdarzeń rzadkich. Z~kolei przy małej próbie i~wielu cechach szczególnego znaczenia nabiera ocena nadmiernego dopasowania. Porównanie metod zostanie zatem uzupełnione analizą ważenia klas i~regularyzacji.

\subsection{Cel i zakres pracy}
\label{sec:cel-zakres}

Celem pracy jest zbadanie przydatności znanych z~literatury metod głębokiego uczenia do klasyfikacji danych tabelarycznych. Ocena będzie oparta na porównaniu z~wybranymi algorytmami konwencjonalnymi oraz na badaniach mechanizmów uwrażliwiania na rzadkie klasy i~ograniczania nadmiernego dopasowania. Praca ma charakter badawczy; nie zakłada konieczności skonstruowania nowej architektury ani wykazania przewagi określonej rodziny metod.

Przyjęto siedem metod neuronowych: MLP, TabNet, FT-Transformer, RealMLP, ModernNCA, TabM i~TabPFN-3.5. Cztery odniesienia konwencjonalne stanowią regresja logistyczna, KNN, XGBoost i~CatBoost. Wybrane metody różnią się architekturą, reprezentacją danych i~sposobem uczenia. RealMLP pozwala uwzględnić znaczenie ustawień całej procedury uczenia, ModernNCA --- uczonej reprezentacji sąsiedztwa, a~TabPFN-3.5, jako przyklad modelu fundamentalnego --- pretrenowania na innych zadaniach.

Badanie obejmie dwanaście zbiorów: Adult, Covertype, Credit Card Fraud, HIGGS w~wybranej wersji podzbioru, Amazon Employee Access, Santander Customer Satisfaction, Online Shoppers Purchasing Intention, Diabetes 130-US Hospitals, Helena, Statlog German Credit, Ionosphere i~Bioresponse. Zachowano zadania klasyfikacji binarnej i~wieloklasowej. Zbiory dobrano celowo, aby uwzględnić zróżnicowane liczności, typy cech i~rozkłady klas; nie stanowią one losowej próby wszystkich zastosowań uczenia maszynowego.

Porównanie główne będzie uzupełnione dwoma badaniami kontrolowanymi. Pierwsze dotyczy ważenia klas przy niezmienionych pozostałych hiperparametrach. Drugie obejmuje usuwanie wybranych składników regularyzacji MLP przy różnych liczbach przykładów treningowych. Aktywne uczenie pozostaje kierunkiem możliwego rozszerzenia, a~nie częścią wymaganą do udzielenia odpowiedzi na przyjęte pytania.

\subsection{Pytania badawcze}
\label{sec:pytania-badawcze}

Przyjęty cel uszczegółowiono w~pięciu pytaniach. Każde odnosi się do zadeklarowanych danych i~warunków oceny, a~nie do wszystkich możliwych modeli i~zastosowań.

\begin{enumerate}
    \item \textbf{PB1:} Jaką jakość predykcji osiągają wybrane metody neuronowe i~konwencjonalne w~dwunastu zadaniach klasyfikacji? Podstawą odpowiedzi będą wyniki poszczególnych zadań i~różnice między procedurami ocenianymi na wspólnych populacjach testowych.
    \item \textbf{PB2:} Jak strojenie hiperparametrów i~łączenie konfiguracji zmieniają jakość oraz koszt? Porównanie obejmie konfigurację domyślną, konfigurację wybraną na walidacji oraz zespół konfiguracji. Warunki te pozostaną rozdzielone w~raportowaniu.
    \item \textbf{PB3:} Jak ważenie klas wpływa na rozpoznawanie zadeklarowanej rzadkiej klasy? Oceniane będą jakość uporządkowania przykładów oraz precyzja i~czułość przy progach wyznaczonych bez dostępu do etykiet testowych.
    \item \textbf{PB4:} Jak dropout i~zanik wag wpływają na generalizację MLP przy różnej liczbie przykładów treningowych? Odpowiedź wymaga porównań z~ustaloną recepturą odniesienia, wspólnych podzbiorów i~kontroli pozostałych ustawień.
    \item \textbf{PB5:} Jaki kompromis jakości, czasu i~pamięci uzyskują badane procedury na zadeklarowanym sprzęcie? Koszt znalezienia rozwiązania zostanie oddzielony od kosztu wykorzystania wybranych modeli i~wykonania predykcji.
\end{enumerate}

Ocenie będą podlegały zarówno poprawa, jak i~pogorszenie jakości oraz przypadki, w~których metoda nie ukończy obliczeń. Brak statystycznego potwierdzenia różnicy nie będzie interpretowany jako równoważność metod.

\subsection{Wkład własny i granice wnioskowania}
\label{sec:wklad}

Wkład własny obejmuje syntezę literatury, uzasadnienie doboru metod i~zadań, projekt porównania oraz integrację odtwarzalnego środowiska eksperymentalnego. Zaimplementowano główną ścieżkę uruchamiania metod, zapisu predykcji i~odtwarzania raportów. Z~uwagi na badawczy charakter pracy, integracja w miarę możliwości korzysta z~istniejących bibliotek i~implementacji; nie oznacza autorstwa oryginalnych algorytmów. Jej poprawność należy rozpatrywać oddzielnie od empirycznej skuteczności metod.

Planowany wkład empiryczny obejmuje porównanie jakości i~kosztu oraz oba badania mechanizmów. Na etapie przygotowania tego tekstu końcowa kampania nie została wykonana. Dostępne sprawdzenia infrastruktury nie są przedstawiane jako wyniki naukowego porównania. Wnioski dotyczące skuteczności zostaną sformułowane dopiero na podstawie zidentyfikowanych artefaktów eksperymentów.

Porównanie kompletnych rozwiązań nie pozwala samo w~sobie przypisać różnicy jakości pojedynczej warstwie, reprezentacji lub regularyzatorowi. Taki wniosek wymaga odpowiednio dopasowanej interwencji. Podobnie jeden podział zewnętrzny nie opisuje zmienności między wszystkimi możliwymi podziałami, a~foldy walidacji wewnętrznej nie zwiększają liczby niezależnych zbiorów w~porównaniu między zadaniami~\cite{demsar2006comparisons}.

\subsection{Organizacja pracy}
\label{sec:organizacja}

Rozdział~\ref{chap:stan-badan} przedstawia właściwości danych tabelarycznych, trudne reżimy klasyfikacji oraz stan badań porównawczych. Wyjaśnia, jakie informacje wynikają z~benchmarków, a~jakie wymagają badań kontrolowanych. W~rozdziale~\ref{chap:metody} opisano mechanizmy wybranych metod i~uzasadniono ich role w~porównaniu.

Rozdział~\ref{chap:metodyka} opisuje podziały danych, strojenie, metryki i~projekt badań kontrolowanych. Rozdział~\ref{chap:srodowisko} przedstawia implementację środowiska eksperymentalnego. Rozdziały~\ref{chap:wyniki} i~\ref{chap:mechanizmy} przygotowano jako szkielety prezentacji wyników; zostaną uzupełnione po wykonaniu eksperymentów. Planowana dyskusja końcowa odniesie wyniki do literatury i~ograniczeń badania.
\clearpage
\section{Specyfika danych tabelarycznych i stan badań}
\label{chap:stan-badan}

\subsection{Definicja zadania i właściwości danych}
\label{sec:definicja-tabel}

Rozpatrywane zadanie opisuje zbiór $\mathcal{D}=\{(x_i,y_i)\}_{i=1}^{n}$, gdzie $x_i\in\mathcal{X}_1\times\cdots\times\mathcal{X}_d$ jest obserwacją o~ustalonym schemacie cech, a~$y_i\in\{1,\ldots,K\}$. Klasyfikator zwraca wektor prawdopodobieństw $p_\theta(y\mid x)$. Oznaczenia $n$, $d$ i~$K$ odnoszą się odpowiednio do liczby obserwacji, cech i~klas; wymiar reprezentacji po kodowaniu może różnić się od~$d$.

Specyfika tabel nie sprowadza się do współwystępowania zmiennych numerycznych i~kategorycznych. Poszczególne kolumny mają odrębne znaczenie, rozkłady i~mechanizmy powstawania braków, a~wybór kodowania określa geometrię wejścia dostępną modelowi~\cite{borisov2022survey}. Traktowanie kodu kategorii jako wielkości porządkowej wprowadza dodatkowe założenie o~relacji między wartościami. W~badaniu znaczenie tego założenia ujawnia się szczególnie w~zestawieniu metod odległościowych, sieci z~osadzeniami i~drzew z~natywną obsługą kategorii.

Odrębną decyzją jest jednostka predykcji. W~Diabetes wiersz opisuje pobyt, podczas gdy rozłączność treningu i~testu ma dotyczyć pacjentów. Identyfikator osoby określa zatem konstrukcję oceny, a~nie reprezentację predykcyjną. Takie rozróżnienie między cechą a~metadaną jest częścią definicji problemu. Tabela~\ref{tab:role-kolumn} zestawia jego konsekwencje dla przygotowania danych.

\begin{table}[!htbp]
\centering
\caption{Typy i role kolumn istotne w projektowaniu klasyfikacji. Zestawienie pojęciowe, nie opis konkretnego zbioru.}
\label{tab:role-kolumn}
\small
\begin{tabularx}{\textwidth}{@{}P{0.22\textwidth}YY@{}}
\toprule
Typ lub rola & Decyzja reprezentacyjna & Konsekwencja dla oceny \\
\midrule
Cecha numeryczna & Skala, transformacja rozkładu, osadzenie & Porównanie obejmuje geometrię wejścia i model. \\
Cecha nominalna & Kodowanie, rzadkie i nieznane poziomy & Słownik i statystyki kategorii zależą od treningu. \\
Cecha porządkowa & Wykorzystanie kolejności lub osobnych poziomów & Odstępy między kodami są dodatkowym założeniem. \\
Identyfikator grupy & Metadana podziału, nie predyktor & Jednostka rozłączności może obejmować wiele wierszy. \\
Czas obserwacji & Moment dostępności informacji & Losowy podział nie identyfikuje jakości przyszłej predykcji. \\
Etykieta docelowa & Semantyka zdarzenia i mapowanie klas & Kod dodatni nie musi oznaczać zdarzenia rzadkiego. \\
\bottomrule
\end{tabularx}
\end{table}

\subsection{Reprezentacja danych i założenia indukcyjne}
\label{sec:reprezentacja-zalozenia}

\subsubsection{Znaczenie reprezentacji kolumn}

W~danych tabelarycznych naturalny układ współrzędnych niesie informację: poszczególne osie odpowiadają określonym pomiarom lub kategoriom. Przestawienie kolumn razem z~ich definicjami jest zmianą zapisu, natomiast obrót mieszający cechy numeryczne zmienia sposób udostępnienia ich znaczenia modelowi. Nie ma podstaw, aby te dwie operacje traktować jako równoważne. Podobnie sąsiedztwo kolumn nie uzasadnia lokalności ani translacyjnego współdzielenia parametrów znanego z~przetwarzania obrazów~\cite{borisov2022survey}.

Kodowanie zero-jedynkowe, osadzenia kategorii i~transformacje cech numerycznych nie są neutralnymi operacjami wejściowymi. Pierwsze zachowuje rozróżnialność kategorii kosztem szerokości macierzy, drugie pozwala uczyć ich relacje, lecz wymaga oszacowania dodatkowych parametrów. Transformacja skali zmienia odległości wykorzystywane przez KNN i~warunki optymalizacji sieci. Z~tego powodu w~porównaniu głównym ocenie podlegają całe potoki, z~przetwarzaniem właściwym danej metodzie. Wymuszenie identycznego kodowania oznaczałoby ocenę innego pytania: jakości modeli przy narzuconej reprezentacji.

\subsubsection{Informacyjność cech i nieregularne zależności}

Grinsztajn i~współautorzy~\cite{grinsztajn2022trees} powiązali przewagę drzew w~badanym reżimie z~trzema własnościami: odpornością na cechy nieinformatywne, zdolnością odtwarzania nieregularnych zależności oraz wykorzystaniem orientacji danych. W~ich interwencjach usuwanie mało informatywnych cech zmniejszało różnicę między sieciami typu MLP a~pozostałymi modelami, natomiast dodawanie cech losowych ją zwiększało. Obrót numerycznego wejścia zmieniał względną kolejność metod, mimo zachowania informacji zawartej w~obserwacjach. Wyniki dotyczą konkretnych architektur i~przekształceń, ale pokazują, dlaczego sama liczba parametrów nie wyjaśnia zdolności uczenia na tabelach.

Z~perspektywy tej pracy prowadzi to do rozdzielenia dwóch problemów: elastyczności reprezentacji i~efektywnego wykorzystania dostępnej próby. Szeroki zbiór nie musi być trudny, jeżeli sygnał jest silny, a~użyteczna zależność prosta. Z~kolei zwiększanie pojemności sieci nie gwarantuje poprawy, gdy dodatkowe stopnie swobody służą dopasowaniu cech nieinformatywnych. Bioresponse stanowi przypadek wysokiej wymiarowości względem liczności treningu, lecz bez kontrolowanych rotacji i~dodawania cech nie pozwala przypisać różnic jakości samemu układowi współrzędnych. Takie rozstrzygnięcie wymagałoby osobnej interwencji.

\subsubsection{Braki i informacje niedostępne przy predykcji}

Brak wartości nie jest tym samym co zero ani nowa kategoria występująca po raz pierwszy podczas predykcji. Procedura musi określać, jak te sytuacje reprezentuje. Jeżeli imputacja lub kodowanie uczą się statystyk, ich dopasowanie powinno respektować granicę między treningiem i~oceną. Włączenie etykiety do kodowania kategorii wymaga szczególnej ostrożności, ponieważ statystyka może ujawnić informację, której model nie otrzyma dla nowego przypadku~\cite{prokhorenkova2018catboost}.

Poprawna obsługa techniczna braków nie wystarcza do stwierdzenia poprawności zadania. Cecha zarejestrowana dopiero po wystąpieniu przewidywanego zdarzenia może zapewniać wysoką jakość bez użyteczności w~zadanym momencie predykcji. Dostępność informacji i~role cech określają dopuszczalną interpretację predykcji; nie wynikają z~samej jakości modelu.

\subsection{Trudności w zadaniach klasyfikacji}
\label{sec:trudne-rezimy}

\subsubsection{Mała liczba przykładów i wysoka wymiarowość}

Dla oceny generalizacji istotna jest liczba przykładów dostępnych w~konkretnym dopasowaniu, a~nie całkowity rozmiar pliku. Po wydzieleniu testu, walidacji i~podzbioru treningowego relacja wymiaru reprezentacji do liczby obserwacji może być znacznie mniej korzystna niż sugeruje opis zbioru. Dodatkowo rzadkie kategorie i~klasy mają własne, mniejsze liczności efektywne. Przy tej samej wartości $n$ dwa zadania mogą zatem dostarczać bardzo różnej ilości informacji o~parametrach decydujących o~predykcji.

Badanie Kadry i~współautorów~\cite{kadra2021regularization} na 40 zbiorach uzasadnia uwzględnienie wspólnie dobieranych technik regularyzacji jako istotnego składnika porównania sieci. Taki wynik nie rozstrzyga jednak, czy dany składnik pomaga przy ustalonych pozostałych ustawieniach. W~badaniu własnym przyjęto właśnie tę drugą perspektywę: usunięcie dropoutu albo zaniku wag z~jednej procedury MLP, przy zmiennej liczności treningu. Przedmiotem oceny jest efekt warunkowy względem architektury, optymalizatora i~reguły zatrzymania, a~nie najlepszy osiągalny wynik każdej osobno strojonej rodziny.

Zmiana udziału treningu przy stałej liczbie epok może zmieniać również liczbę aktualizacji parametrów. Jest to istotne dla interpretacji zarówno optymalizacji, jak i~skumulowanego zaniku wag. Dlatego krzywe jakości względem liczności są interpretowane łącznie z~przebiegiem uczenia i~wybraną epoką. Stałe walidacja i~test oddzielają zmianę treningu od zmiany ocenianych populacji. Taka zależność stanowi krzywą wrażliwości na liczność treningu, przy zachowanym budżecie etykiet walidacyjnych.

\subsubsection{Niezrównoważenie i rozpoznawanie rzadkiej klasy}

W~klasyfikacji niezrównoważonej trzeba rozdzielić trzy obiekty oceny: uporządkowanie obserwacji, wartości prawdopodobieństw i~decyzje przy zadanym progu. ROC AUC opisuje pierwszy z~nich, ale nie określa udziału prawdziwych zdarzeń wśród alarmów. Przy częstości klasy rzadkiej $\pi$ precyzja dla ustalonego progu spełnia tożsamość
\begin{equation}
\operatorname{Precision}=\frac{\pi\operatorname{TPR}}
{\pi\operatorname{TPR}+(1-\pi)\operatorname{FPR}},
\label{eq:precyzja-czestosc}
\end{equation}
gdzie TPR oznacza czułość, a~FPR odsetek fałszywych alarmów w~klasie pozostałej. Zależność wynika bezpośrednio z~definicji prawdopodobieństwa warunkowego: przy dodatnich, niezmienionych TPR i~FPR spadek $\pi$ obniża precyzję. Uzasadnia to analizę precision--recall obok ROC, zgodnie z~argumentacją Saito i~Rehmsmeiera~\cite{saito2015pr}. W~badaniu miarą oceny jest average precision (AP), liczona jako suma przyrostów czułości ważonych precyzją, a~nie jako trapezowa całka krzywej.

Rozróżnienie ważenia i~progowania można doprecyzować analitycznie. Niech $\eta(x)=P(Y=1\mid X=x)$ oznacza prawdopodobieństwo rzadkiego zdarzenia, przy czym kod~1 przyjęto wyłącznie na potrzeby poniższego wyprowadzenia. Dla dodatnich, stałych wag klas $w_1,w_0$ warunkowe ryzyko ważonej straty logarytmicznej ma postać
\begin{equation}
R_w(q\mid x)=-w_1\eta(x)\log q-w_0[1-\eta(x)]\log(1-q).
\label{eq:ryzyko-wazone}
\end{equation}
Wyzerowanie pochodnej po $q$ daje, dla $0<\eta(x)<1$,
\begin{equation}
q_w^*(x)=\frac{w_1\eta(x)}{w_1\eta(x)+w_0[1-\eta(x)]},
\qquad
\operatorname{logit}q_w^*(x)=\operatorname{logit}\eta(x)+\log\frac{w_1}{w_0}.
\label{eq:optimum-wazenia}
\end{equation}
Powyższe wyprowadzenie dotyczy nieograniczonej klasy funkcji i~ryzyka populacyjnego, nie rzeczywistego przebiegu optymalizacji. W~tych warunkach ważenie zmienia skalę prawdopodobieństw, lecz zachowuje uporządkowanie: transformacja jest ściśle rosnąca. Próg $q_w^*\geq 1/2$ odpowiada progowi $\eta\geq w_0/(w_0+w_1)$. Sam wzrost czułości po ważeniu nie dowodzi więc lepszej reprezentacji ani lepszego rankingu.

W~skończonej próbie, przy ograniczonej pojemności i~regularyzacji, ważenie może zmienić również wyuczone uporządkowanie. Dlatego badanie~A obejmuje równolegle AP, ROC AUC, nieważoną stratę logarytmiczną oraz decyzje przy progach wybranych na walidacji. Pozwala to odróżnić poprawę rankingu od przesunięcia punktu pracy i~zmiany jakości prawdopodobieństw. Zdarzenie rzadkie jest identyfikowane semantycznie, niezależnie od kodu nadanego mu przez procedurę przygotowania danych.

\subsubsection{Wieloklasowość i jakość prawdopodobieństw}

W~zadaniu wieloklasowym strata logarytmiczna ocenia cały rozkład predykcyjny poprzez prawdopodobieństwo przypisane etykiecie rzeczywistej. Dla ustalonego $x$, prawdziwego rozkładu $\eta_c=P(Y=c\mid x)$ i~predykcji $p_c$ jej wartość oczekiwaną można zapisać jako
\begin{equation}
-\sum_{c=1}^{K}\eta_c\log p_c
=H(\eta)+D_{\mathrm{KL}}(\eta\Vert p).
\label{eq:logloss-kl}
\end{equation}
Dekompozycja wynika z~dodania i~odjęcia $\sum_c\eta_c\log\eta_c$. Przy dostępie do prawdziwego rozkładu minimum osiąga $p=\eta$; wysoka pewność błędnej klasy jest karana nawet wtedy, gdy liczba poprawnych decyzji pozostaje bez zmian. Niższa strata na skończonym teście nie identyfikuje jednak osobno poprawy kalibracji: zależy również od zdolności rozróżniania przypadków.

W~przyjętym protokole strata logarytmiczna jest główną miarą zadań wieloklasowych, zgodnie z~adaptacją TabArena~\cite{erickson2025tabarena}. Covertype, Diabetes i~Helena reprezentują różne liczby klas. W~Helenie istotna jest nie tylko obsługa stu wyjść, lecz także rozkład liczności klas i~jakość oszacowania prawdopodobieństw przy ograniczonej liczbie przykładów niektórych etykiet. Wartości straty są porównywane między metodami w~obrębie zadania; ich bezpośrednie uśrednienie między zadaniami mieszałoby różne poziomy trudności i~entropii etykiet.

\subsubsection{Zależności między obserwacjami i zmiana rozkładu}

Powtarzające się obserwacje tej samej osoby mogą naruszać interpretację losowego podziału jako oceny nowych osób. Podobnie losowe przemieszanie danych z~różnych okresów nie musi odpowiadać przewidywaniu przyszłości. Rubachev i~współautorzy~\cite{rubachev2025tabred} wskazali znaczenie podziałów czasowych w~TabReD, obejmującym osiem przemysłowych zbiorów z~rozbudowanymi cechami. Zmiana sposobu podziału wpływała na względną ocenę metod.

W~tej pracy podział Diabetes zachowuje rozłączność pacjentów. Zakres badania nie obejmuje osobnego porównania czasowego, dlatego nie uzasadnia wniosków o~odporności na dryf. Dla pozostałych źródeł interpretacja wyniku również zależy od tego, czy podział odzwierciedla właściwą populację predykcji.

\subsection{Rozwój podejść neuronowych}
\label{sec:rozwoj-podejsc}

Rozwiązania neuronowe można porządkować według tego, na którym poziomie ingerują w~procedurę predykcji. Zmiana reprezentacji przygotowuje wejście dla modelu, zmiana architektury określa sposób łączenia informacji, a~procedura uczenia wpływa na odnajdywane parametry. Podział ten jest użyteczny, lecz nie jest rozłączny: jedna metoda może łączyć kilka takich elementów~\cite{borisov2022survey}.

Przyjęty dobór obejmuje także podejście oparte na sąsiedztwie i~model wykorzystujący wiedzę nabytą przed otrzymaniem badanego zadania. Ich obecność rozszerza porównanie poza klasyczny schemat trenowania niezależnej sieci od początku na każdym zbiorze. Podział na modele neuronowe i~konwencjonalne pozostaje przydatny organizacyjnie, ale nie wyczerpuje różnic między procedurami.

Rysunek~\ref{fig:mapa-podejsc} porządkuje wybrane metody według ich roli w~porównaniu. Wyjaśnienie mechanizmów znajduje się w~rozdziale~\ref{chap:metody}.

\plannedfigure
    {figures/ch02/mapa-podejsc.pdf}
    {Role wybranych podejść w~badaniu. Opracowanie własne.}
    {fig:mapa-podejsc}
    {\begin{tabular}{ll}
Odniesienie gęste & MLP \\
Selekcja cech i~uwaga & TabNet, FT-Transformer \\
Procedura uczenia & RealMLP \\
Uczone sąsiedztwo & ModernNCA \\
Współdzielenie parametrów zespołu & TabM \\
Pretrenowanie i~kontekst & TabPFN-3.5 \\
Odniesienia konwencjonalne & Regresja logistyczna, KNN, XGBoost, CatBoost
\end{tabular}}

\subsection{Badania porównawcze}
\label{sec:benchmarki}

Przegląd ma charakter narracyjny i~obejmuje publikacje uwzględnione przy doborze metod do 27 października 2026 roku. Dla preprintów wskazano wykorzystane wersje, ponieważ kolejne rewizje mogą różnić się zakresem badań. Przegląd nie był systematycznym przeszukaniem całej literatury dziedziny.

\subsubsection{Od porównania architektur do analizy warunków}

Gorishniy i~współautorzy~\cite{gorishniy2021revisiting} porównali architektury na jedenastu zadaniach klasyfikacji i~regresji we wspólnej procedurze uczenia i~strojenia. Istotnym wynikiem było wykazanie konkurencyjności prostego ResNet oraz silnej pozycji FT-Transformer wśród badanych sieci. Porównanie z~drzewami nie wskazało rozwiązania wygrywającego we wszystkich zadaniach. Dla projektu tej pracy oznacza to konieczność uwzględnienia mocnego odniesienia gęstego: specjalizacja architektury wymaga wykazania korzyści względem właściwie przygotowanej sieci bazowej.

Grinsztajn i~współautorzy~\cite{grinsztajn2022trees} uzyskali przewagę drzew na benchmarku 45 zbiorów, utrzymującą się wraz ze zwiększaniem liczby prób strojenia. Badanie obejmowało jednak określony reżim: ograniczoną liczność treningu, usunięte braki, równoważenie klasyfikacji oraz ograniczenia wymiarowości i~kategorii. Wynik stanowi argument za zachowaniem silnych odniesień drzewiastych, lecz nie rozstrzyga o~naturalnie rzadkich klasach ani o~wszystkich późniejszych procedurach neuronowych.

McElfresh i~współautorzy~\cite{mcelfresh2024tabzilla} rozszerzyli perspektywę do 19 algorytmów i~176 zbiorów klasyfikacji. W~ich analizie na około jednej trzeciej zadań przyrost jakości po lekkim strojeniu CatBoost był większy niż różnica między najlepszym domyślnym modelem drzewiastym i~neuronowym. Znaczenie ma tu konkretna konstrukcja porównania, a~nie ogólne pierwszeństwo strojenia przed doborem algorytmu. Dla niniejszej pracy wynik uzasadnia oddzielną ocenę ustawień domyślnych i~strojonych. Wczesny TabPFN był przy tym oceniany w~odmiennych warunkach i~nie stanowi dowodu skuteczności wybranej tutaj wersji~3.5.

\subsubsection{Procedury uczenia i rozwój benchmarków}

Holzmüller i~współautorzy~\cite{holzmueller2024realmlp} opracowali ustawienia RealMLP na 118 zbiorach, a~następnie ocenili ich przenoszenie na 90 zbiorów meta-test. Uzyskali konkurencyjność rozwiniętej procedury gęstej względem drzew, lecz obraz zależał od kolekcji i~metryki: na benchmarku Grinsztajna RealMLP wypadał słabiej od CatBoost w~klasyfikacji. Ważnym ustaleniem jest zatem skuteczność przenoszenia ustawień między zadaniami, a~nie bezwarunkowe rozstrzygnięcie na korzyść MLP. Konfiguracja domyślna może zawierać rezultat kosztownego wcześniejszego rozwoju, nawet jeżeli nie jest strojona na aktualnym zbiorze.

Wersja~4 pracy Ye i~współautorów~\cite{ye2025talent} obejmuje 40 metod na 300 bazowych zbiorach oraz dodatkowe kolekcje. W~zestawieniach wysokie pozycje zajmują modele pretrenowane, a~wśród metod uczonych zadaniowo wyróżniają się rozwinięte MLP i~ModernNCA. Starsze specjalizacje, w~tym TabNet, wypadają słabiej od czołowych zespołów. Autorzy odnotowują także korzyści zespołowania, między innymi dla TabM. To uzasadnia reprezentowanie kilku kierunków rozwoju sieci zamiast utożsamiania całej rodziny z~jednym MLP. Główny protokół stosuje podział 64/16/20 i~100 prób strojenia; agregaty obejmujące regresję pozostają odrębne od wniosków o~klasyfikacji.

TabArena~\cite{erickson2025tabarena} przenosi punkt ciężkości z~jednorazowego rankingu na wersjonowany system oceny danych, modeli i~procedur wyszukiwania. Początkowa kolekcja obejmuje 51 zbiorów. Dla niniejszego projektu najistotniejsze jest rozdzielenie konfiguracji domyślnej, wyboru po strojeniu i~zespołu konfiguracji. Pozwala ono ocenić, czy poprawa wynika z~dodatkowego wyszukiwania, czy z~wykorzystania różnorodności znalezionych rozwiązań. Przyjęte dwanaście zadań i~własne wyjątki walidacyjne tworzą adaptację tej organizacji eksperymentu, nie reprodukcję oficjalnego rankingu.

\begin{table}[tbp]
\centering
\caption{Wybrane badania porównawcze i ich znaczenie dla projektu pracy. Liczby dotyczą wskazanych publikacji lub wersji, a nie jednego wspólnego benchmarku. Opracowanie własne na podstawie cytowanych źródeł.}
\label{tab:benchmarki}
\footnotesize
\begin{tabularx}{\textwidth}{@{}P{0.25\textwidth}P{0.22\textwidth}Y@{}}
\toprule
Praca & Zakres & Ograniczenia porównania \\
\midrule
RTDL~\cite{gorishniy2021revisiting} & 11 zadań & Klasyfikacja i regresja; starsze modele i inny protokół. \\
Grinsztajn i in.~\cite{grinsztajn2022trees} & 45 zbiorów & Kontrolowany reżim nie zachowuje wszystkich trudności danych źródłowych. \\
TabZilla~\cite{mcelfresh2024tabzilla} & 176 zbiorów klasyfikacji & Poszczególne analizy obejmują różne podzbiory zadań i ustawienia modeli. \\
RealMLP~\cite{holzmueller2024realmlp} & 118 meta-train, 90 meta-test & Rozwój procedury jest odrębny od oceny i kosztów zadaniowych. \\
TALENT v4~\cite{ye2025talent} & 300 zbiorów bazowych i rozszerzenia & Mieszane typy zadań i inna walidacja. \\
TabArena~\cite{erickson2025tabarena} & 51 zbiorów początkowych & Zespoły, wersje i zasoby są częścią ocenianej procedury. \\
TabReD~\cite{rubachev2025tabred} & 8 zbiorów przemysłowych & Podziały czasowe nie odpowiadają standardowemu IID. \\
\bottomrule
\end{tabularx}
\end{table}

\subsection{Zgodności, rozbieżności i ograniczenia syntezy}
\label{sec:synteza-literatury}

Z~zestawienia wynika przede wszystkim, że etykieta rodziny modelu jest zbyt zgrubna, aby wyjaśniać postęp. Prace RTDL i~RealMLP podważają ocenę nowych architektur względem słabo skonfigurowanej sieci gęstej; TALENT pokazuje znaczenie rozwiniętych procedur i~zespołów. Wspólnym wnioskiem przyjętym w~tej pracy jest konieczność porównania konkretnych konfiguracji wraz z~ich reprezentacją i~kosztem selekcji, a~nie abstrakcyjnych kategorii ,,sieci'' i~,,drzewa''~\cite{gorishniy2021revisiting,holzmueller2024realmlp,ye2025talent}.

Rozbieżność między przewagą drzew w~badaniu Grinsztajna a~konkurencyjnością nowszych sieci nie jest prostym odwróceniem jednego wyniku. Zmieniono zarówno procedury, jak i~populacje zadań, budżety oraz kryteria agregacji. RealMLP dostarcza szczególnie użytecznego przykładu: jego względna pozycja zmienia się między kolekcjami ocenianymi w~tej samej publikacji~\cite{grinsztajn2022trees,holzmueller2024realmlp}. W~niniejszym badaniu uzasadnia to prezentację różnic zadaniowych obok rang. Nie przeprowadzono metaanalizy wyznaczającej wspólną wielkość przewagi z~tych niejednorodnych protokołów.

Pozostaje również różnica między wyjaśnieniem wyniku a~wyborem skutecznej procedury. Benchmark ocenia efekt całego rozwiązania; interwencja usuwa lub zmienia wskazany składnik. Wyniki tych dwóch rodzajów badań odpowiadają na uzupełniające się pytania. Powiązanie rankingu z~metacechami może sugerować hipotezę, lecz dopiero dopasowana kontrola pozwala odnieść różnicę do zadanej zmiany. Właśnie na tej granicy oparto podział własnego badania na porównanie główne i~dwa eksperymenty uzupełniające.

Ważną osią porównywalności jest oddzielenie selekcji od oceny. Cawley i~Talbot~\cite{cawley2010selection} pokazują, że dopasowywanie się do skończonej próby walidacyjnej może obciążać ocenę wybranej procedury. Wspólna walidacja nie daje zatem prawa do wielokrotnego projektowania zmian na podstawie testu. Etykiety testowe nie powinny służyć wyborowi konfiguracji, progu, zespołu ani zakresu prezentowanych porównań.

Drugą osią jest jednostka wnioskowania. Porównanie między zbiorami różni się od analizy powtórzeń na jednym zbiorze~\cite{demsar2006comparisons}. W~niniejszym badaniu foldy i~konfiguracje nie są traktowane jako nowe niezależne zadania. Rangi uzupełniają wyniki zadaniowe, ponieważ pomijają wielkość różnicy. Ponadto brak informacji o~pokryciu danych rozwoju metod nie jest zastępowany założeniem, że wszystkie wybrane zadania są całkowicie nowe dla ich twórców.

\subsection{Uzasadnienie własnego projektu}
\label{sec:uzasadnienie-projektu}

Własne badanie łączy ocenę użyteczności kompletnych procedur z~analizą dwóch mechanizmów istotnych dla trudnych zadań klasyfikacji. Porównanie główne odpowie, jaką jakość i~koszt uzyskuje się przy określonym sposobie selekcji modelu. Nie wyjaśni jednak, czy ewentualny wzrost wykrywalności rzadkiej klasy wynika z~lepszego rankingu, ani czy zachowanie MLP przy małej próbie zależy od danego składnika regularyzacji. Temu służą ważenie klas oraz dopasowane ablacje przy zagnieżdżonych podzbiorach treningowych.

Wartość tego układu polega na wspólnej interpretacji jakości, kosztu i~reakcji na kontrolowaną zmianę. Nie zakłada się, że każdy wybrany model jest skuteczny w~każdym reżimie. Możliwym wynikiem jest zarówno korzyść mechanizmu, jak i~jej brak lub odwrócenie po zmniejszeniu treningu. Tabela~\ref{tab:problem-dowod} określa, które pytania mają bezpośrednie pokrycie w~projekcie eksperymentu, a~które stanowią kontekst interpretacji.

\begin{table}[tbp]
\centering
\caption{Powiązanie właściwości danych z~zakresem badania.}
\label{tab:problem-dowod}
\small
\begin{tabularx}{\textwidth}{@{}P{0.26\textwidth}YY@{}}
\toprule
Problem & Zakres porównania & Granica interpretacji \\
\midrule
Jakość i koszt potoków & Wspólne testy, D/T/T+E & Bez izolacji pojedynczej architektury. \\
Rzadkie klasy & Ważenie i kontrola progów & Efekt przy ustalonych pozostałych ustawieniach. \\
Mało danych i regularyzacja & Zagnieżdżone podzbiory, ablacje MLP & Bez pełnej analizy interakcji mechanizmów. \\
Wymiarowość i typy cech & Zróżnicowane naturalne zadania & Związek opisowy, nie kontrolowany efekt jednej cechy. \\
Szum, rotacje, dryf & Omówienie literatury i audyt & Brak własnego testu odporności w tych reżimach. \\
\bottomrule
\end{tabularx}
\end{table}

Badanie główne pozwala porównać jakość i~koszt wybranych metod w~jednym protokole. Badania uzupełniające zawężają analizę do wpływu ważenia klas oraz dwóch składników regularyzacji MLP. Mechanizmy tych metod i~uzasadnienie ich wyboru przedstawiono w~kolejnym rozdziale.

\clearpage
\section{Wybrane metody klasyfikacji}
\label{chap:metody}

\subsection{Zasady doboru i wspólna terminologia}
\label{sec:dobor-metod}

Metody wybrano ze względu na ich role w~porównaniu, a~nie wyłącznie datę publikacji lub dostępność kodu. Pula obejmuje siedem rozwiązań neuronowych i~cztery konwencjonalne. Pozwala ona porównać liniową klasyfikację, sąsiedztwo, zespoły drzew oraz odmienne sposoby wykorzystania sieci. Nie obejmuje wszystkich rodzin opisanych w~literaturze; brak metody w~macierzy nie oznacza jej nieskuteczności.

Przez \emph{potok uczenia} rozumie się tutaj przygotowanie danych, dopasowanie klasyfikatora i~regułę wyznaczania predykcji. Parametry modelu są wielkościami uzyskiwanymi w~trakcie dopasowania. Hiperparametry określają m.in. strukturę, regularyzację i~przebieg tego procesu, a~ich wybór stanowi osobny poziom selekcji. Konfiguracja domyślna jest wskazanym punktem odniesienia, lecz może zawierać ustawienia opracowane wcześniej na innych zadaniach.

Każda metoda ma zwracać prawdopodobieństwa w~jednakowej, zapisanej kolejności klas. Daje to wspólny interfejs oceny bez wymuszania identycznej reprezentacji wejścia. Narzucenie jednego kodowania wszystkim algorytmom mogłoby usunąć mechanizm, który uzasadnia ich wybór. Z~tego względu porównanie główne dotyczy kompletnych procedur. Tabela~\ref{tab:role-metod} podsumowuje ich role; szczegółowe ustawienia wykonania należą do metodyki i~artefaktów badania.

\begin{table}[tbp]
\centering
\caption{Metody przyjęte do porównania i ich role badawcze. Opracowanie własne; zestawienie nie jest rankingiem.}
\label{tab:role-metod}
\small
\begin{tabularx}{\textwidth}{@{}P{0.22\textwidth}P{0.24\textwidth}Y@{}}
\toprule
Metoda & Mechanizm & Rola w porównaniu \\
\midrule
Regresja logistyczna & Liniowa ocena klas & Odniesienie liniowe w przyjętej reprezentacji cech. \\
KNN & Sąsiedztwo przykładów & Konwencjonalne odniesienie do uczonego sąsiedztwa. \\
XGBoost & Wzmacniane drzewa & Silne odniesienie drzewiaste. \\
CatBoost & Drzewa i obsługa kategorii & Odniesienie zachowujące dedykowane przetwarzanie kategorii. \\
MLP & Warstwy gęste & Bazowa metoda neuronowa. \\
TabNet & Sekwencyjne maski cech & Specjalizacja przez selekcję informacji. \\
FT-Transformer & Tokeny cech i uwaga & Specjalizacja reprezentacji i interakcji kolumn. \\
RealMLP & Rozwinięta receptura MLP & Znaczenie przygotowania danych i uczenia. \\
ModernNCA & Uczona reprezentacja sąsiedztwa & Porównanie z klasycznym KNN. \\
TabM & Zespół ze współdzieleniem parametrów & Efektywne łączenie predykcji wewnątrz architektury. \\
TabPFN-3.5 & Pretrenowanie i kontekst zadania & Zewnętrznie nabyte wagi przy stałej adaptacji. \\
\bottomrule
\end{tabularx}
\end{table}

\subsection{Metody konwencjonalne}
\label{sec:metody-konwencjonalne}

\subsubsection{Regresja logistyczna}
\label{sec:regresja-logistyczna}

Regresja logistyczna wiąże liniową kombinację cech z~prawdopodobieństwem klasy. W~przypadku binarnym stosowana jest funkcja logistyczna, natomiast uogólnienie wieloklasowe pozwala uzyskać rozkład po etykietach. Regularyzacja ogranicza wartości współczynników. Takie własności opisuje również dokumentacja implementacji scikit-learn~\cite{sklearnLogistic}.

Model stanowi odniesienie pozwalające sprawdzić, czy do uzyskania dobrej jakości potrzebna jest bardziej złożona zależność niż liniowa w~przyjętej reprezentacji. Nie oznacza to liniowości względem surowego zapisu każdej kolumny: kodowanie kategorii zmienia wejście. W~przyjętym potoku nominalne wartości są kodowane zero-jedynkowo, a~przekształcenia dopasowywane na odpowiedniej części treningowej. Dobra jakość tego odniesienia ogranicza uzasadnienie kosztowniejszych metod; słaba jakość sama w~sobie nie dowodzi przewagi określonej architektury.

\subsubsection{Klasyfikacja najbliższych sąsiadów}
\label{sec:knn}

Reguła najbliższego sąsiada przypisuje nowemu przypadkowi klasę na podstawie podobnych oznakowanych obserwacji~\cite{cover1967nearest}. KNN rozszerza tę zasadę na wybrany zbiór sąsiadów i~agreguje ich etykiety, ewentualnie z~wagami zależnymi od odległości~\cite{sklearnNeighbors}. Kluczowe znaczenie mają liczba sąsiadów, metryka i~reprezentacja cech. Duża skala jednej kolumny może zdominować odległość, a~dowolne kody kategorii mogą narzucić nieuzasadnione sąsiedztwo.

W~badaniu nominalne kategorie KNN mają ustalone kodowanie zero-jedynkowe. Należy uwzględniać pamięć banku odniesienia i~koszt predykcji, a~nie tylko przygotowania klasyfikatora. Natywna walidacja LOO dla zadań IID wyłącza etykietę ocenianej obserwacji z~jej sąsiedztwa, lecz dopasowuje preprocessing na całej części deweloperskiej. Jest to jawny wyjątek walidacyjny, odmienny od dokładnego ponownego dopasowania wszystkich przekształceń po usunięciu każdego wiersza. Dla Diabetes zastosowane zostaną foldy grupowe.

\subsubsection{XGBoost}
\label{sec:xgboost}

XGBoost jest systemem wzmacniania drzew, w~którym predyktor powstaje przez kolejne dodawanie funkcji drzewiastych. Nowe drzewa ograniczają stratę aktualnego zespołu, a~kryterium uwzględnia regularyzację jego złożoności. Chen i~Guestrin~\cite{chen2016xgboost} opisują również mechanizmy usprawniające uczenie na danych rzadkich i~dużych kolekcjach.

Podziały drzew umożliwiają modelowanie progów i~interakcji bez przyjmowania globalnej liniowości. Nie usuwa to potrzeby doboru głębokości, tempa uczenia, liczby drzew i~regularyzacji. W~pracy XGBoost stanowi silne odniesienie konwencjonalne, oceniane z~odpowiednią funkcją klasyfikacji i~zadeklarowaną reprezentacją. Własności współczesnej biblioteki nie są utożsamiane z~każdym szczegółem publikacji z~2016 roku; wersja wykonania jest odrębnym elementem odtwarzalności.

\subsubsection{CatBoost}
\label{sec:catboost}

CatBoost łączy wzmacnianie drzew z~obsługą cech kategorycznych. Prokhorenkova i~współautorzy~\cite{prokhorenkova2018catboost} opisują statystyki kategorii wyznaczane według porządku obserwacji oraz uporządkowane wzmacnianie. Celem tych mechanizmów jest ograniczanie określonych źródeł przecieku i~przesunięcia predykcji związanych z~użyciem etykiet treningowych.

W~badaniu zostanie zachowana natywna obsługa cech kategorycznych CatBoost. Metoda będzie punktem odniesienia w~porównaniu głównym oraz w~badaniu ważenia klas. Efektywna konfiguracja określi, które warianty mechanizmów opisanych w~publikacji są używane przez bibliotekę.

\subsection{MLP i znaczenie receptury uczenia}
\label{sec:mlp}

Perceptron wielowarstwowy przetwarza wektor wejściowy przez kolejne warstwy gęste z~nieliniowymi aktywacjami. W~klasyfikacji końcowa warstwa wyznacza oceny klas przekształcane w~prawdopodobieństwa. Uczenie parametrów opiera się na minimalizacji straty z~wykorzystaniem gradientów. Prosty MLP jest istotnym odniesieniem, ponieważ pozwala ocenić, czy bardziej wyspecjalizowane mechanizmy dają korzyść względem odpowiednio skonfigurowanej sieci gęstej~\cite{rtdlImplementation}.

Lokalny model bazowy wykorzystuje aktywacje ReLU, dropout i~optymalizator AdamW. Konfiguracja domyślna zawiera dwie warstwy ukryte o~256 i~128 jednostkach, dropout równy~0,1, tempo uczenia~0,001 i~zanik wag~$10^{-4}$. Cechy numeryczne są skalowane, a~kategorie kodowane zero-jedynkowo z~limitem 256 kolumn na cechę. Uczenie trwa najwyżej 100 epok; zachowywany jest stan sieci o~najniższej stracie walidacyjnej. Są to ustawienia własnego modelu bazowego, a~nie dokładne odtworzenie jednej z~opublikowanych procedur MLP. W~wariancie strojonym część tych wartości jest wybierana z~zadeklarowanej przestrzeni hiperparametrów.

Dropout losowo wyłącza jednostki w~trakcie treningu, zmieniając reprezentacje dostępne w~kolejnych aktualizacjach~\cite{srivastava2014dropout}. Podczas zwykłej oceny jest wyłączony; procedura uwzględnia odpowiednie skalowanie. Zanik wag ogranicza parametry w~trakcie optymalizacji. W~AdamW jest oddzielony od aktualizacji wynikającej z~gradientu straty; nie należy utożsamiać go bez zastrzeżeń z~dodaniem kary $L_2$ do straty optymalizowanej Adamem~\cite{loshchilov2019adamw}.

Wpływ regularyzacji zależy od innych ustawień. Zbyt silne ograniczenie może prowadzić do niedopasowania. Niższy błąd treningowy nie musi oznaczać lepszej jakości na zbiorze testowym, a~mniejsza luka generalizacji może wynikać z~pogorszenia dopasowania treningu. Potrzebna jest wspólna analiza jakości na treningu i~na danych niewykorzystanych do aktualizacji parametrów.

W~osobnym badaniu zostanie zastosowana ustalona procedura uczenia MLP oparta na opublikowanym rozwiązaniu oraz dwie ablacje: wyłączenie dropoutu i~wyłączenie zaniku wag. Pozostałe ustawienia będą zachowane. Dobór tej receptury nie będzie oparty na wyniku strojenia z~pełną liczbą etykiet trzech badanych zadań. Wyniki wieloskładnikowych zestawów regularyzacji~\cite{kadra2021regularization} są motywacją takiego rozdzielenia, lecz nie określają automatycznie wartości parametrów własnej kontroli.

\begin{table}[tbp]
\centering
\caption{Dwa zastosowania MLP w pracy. Rozdzielenie modelu bazowego od kontrolowanego badania regularyzacji.}
\label{tab:mlp-role}
\small
\begin{tabularx}{\textwidth}{@{}P{0.24\textwidth}YY@{}}
\toprule
Właściwość & MLP w porównaniu głównym & MLP w badaniu regularyzacji \\
\midrule
Cel & Odniesienie neuronowe & Ocena skutku usunięcia składnika. \\
Dobór ustawień & Domyślne i strojenie wewnętrzne & Receptura ustalona niezależnie. \\
Liczność treningu & Wynik przyjętego podziału & Zagnieżdżone podzbiory treningu. \\
Kontrola zmian & Porównanie pełnych potoków & Jedna ablacja przy stałych pozostałych warunkach. \\
\bottomrule
\end{tabularx}
\end{table}

\subsection{TabNet i FT-Transformer}
\label{sec:tabnet-ft}

\subsubsection{TabNet: sekwencyjna selekcja cech}
\label{sec:tabnet}

TabNet wykorzystuje kolejne kroki przetwarzania, w~których uczone maski określają udział cech dla danej obserwacji. Informacja z~poprzedniego kroku wpływa na następną selekcję, a~częściowe reprezentacje współtworzą wynik. Mechanizm ma kierować zdolność modelowania na wybrane składniki wejścia zamiast przetwarzać je jednakowo w~każdym kroku~\cite{arik2021tabnet}.

Maski mogą służyć opisowi zachowania modelu, lecz nie są dowodem przyczynowego znaczenia cech. W~tej pracy TabNet reprezentuje historycznie ważną specjalizację dla tabel. Nie zaplanowano osobnego badania wierności wyjaśnień ani etapu samonadzorowanego pretrenowania, również opisanego w~oryginalnej publikacji. Wykorzystana implementacja biblioteczna będzie identyfikowana oddzielnie od publikacji metody.

\subsubsection{FT-Transformer: tokeny cech i uwaga}
\label{sec:ft-transformer}

FT-Transformer przekształca każdą cechę w~token o~wspólnym wymiarze. Dla cech numerycznych stosuje parametry właściwe kolumnie, a~dla kategorii --- uczone tablice osadzeń. Zachowuje to tożsamość cechy: jednakowa wartość liczbowa w~dwóch kolumnach nie musi otrzymać tej samej reprezentacji. Warstwy uwagi przetwarzają tokeny cech pojedynczej obserwacji, a~reprezentacja tokenu agregującego służy predykcji~\cite{gorishniy2021revisiting}.

Mechanizm pełnej uwagi wymaga rozpatrywania par tokenów, dlatego szerokość tabeli ma znaczenie dla kosztu. W~porównaniu z~TabNet inny jest sposób selekcji i~łączenia informacji. Różnica całych wyników nie wyizoluje jednak skutku samego rodzaju uwagi. Adapter w~repozytorium korzysta z~implementacji autorów; liczby bloków, wymiar tokenu i~ustawienia uczenia muszą być zapisane jako element wykonanej konfiguracji~\cite{rtdlImplementation}.

Na rysunku~\ref{fig:reprezentacje} przewidziano porównanie przepływu informacji w~MLP, TabNet i~FT-Transformer. Jest to schemat mechanizmów, a~nie ilustracja przewagi jednej metody.

\plannedfigure
    {figures/ch03/przeplyw-reprezentacji.pdf}
    {Przetwarzanie cech w MLP, TabNet i FT-Transformer. Opracowanie własne na podstawie opisów metod~\cite{gorishniy2021revisiting,arik2021tabnet}.}
    {fig:reprezentacje}
    {Trzy równoległe ścieżki od tej samej obserwacji do prawdopodobieństw klas. Pokazać wektor cech, kolejne maski lub tokeny z tożsamością kolumn.}

\subsection{RealMLP: rozwinięta receptura sieci gęstej}
\label{sec:realmlp}

RealMLP jest zestawem rozwiązań obejmujących przetwarzanie wejścia, konstrukcję sieci i~uczenie. W~publikacji opisano m.in. skalowanie odporne z~łagodnym ograniczaniem wartości, osadzenia i~zmodyfikowane ustawienia optymalizacji~\cite{holzmueller2024realmlp}. Nazwa nie powinna być interpretowana jako oznaczenie jednej dodatkowej warstwy dołączonej do dowolnego MLP.

W~badaniu RealMLP odpowiada na pytanie o~jakość rozwiniętego rozwiązania gęstego względem lokalnej sieci i~odniesień konwencjonalnych. Warianty domyślne i~strojone będą wynikać z~przypiętego adaptera i~jego przestrzeni konfiguracji. Konfiguracje do strojenia generuje biblioteka TabArena; lokalny plik YAML nie definiuje dla tej metody osobnej przestrzeni.

Ustawienia opracowane na innych zadaniach są formą wiedzy zewnętrznej. Wynik MLP--RealMLP będzie zatem różnicą kompletnych potoków, obejmującą również przygotowanie danych i~recepturę. Nie posłuży jako jednoczynnikowy dowód działania konkretnego regularyzatora; temu służy osobne badanie MLP.

\subsection{ModernNCA: uczenie reprezentacji sąsiedztwa}
\label{sec:modernnca}

ModernNCA rozwija ideę Neighbourhood Components Analysis, zastępując wyłącznie liniową reprezentację rozwiązaniem wykorzystującym współczesne składniki głębokiego uczenia. Przewidywanie wykorzystuje miękką agregację etykiet sąsiadów w~uczonej przestrzeni. Autorzy stosują również losowe próbkowanie kandydatów sąsiedztwa, aby ograniczać koszt treningu~\cite{ye2025modernnca}.

Sieć uczy się reprezentacji, w~której wyznaczane jest sąsiedztwo. Podczas predykcji nadal potrzebne są oznakowane przykłady odniesienia, dlatego koszt metody obejmuje również ich przechowywanie i~porównywanie z~nowymi obserwacjami.

KNN stanowi bezpośrednie odniesienie, lecz porównanie obejmie różne sposoby uczenia, przetwarzania i~agregacji. Nie wyizoluje wyłącznie korzyści reprezentacji. Kandydaci i~ich etykiety muszą pochodzić z~treningu bieżącego foldu; etykieta walidacyjna lub testowa nie może trafić do banku. W~środowisku użyto implementacji wywodzącej się z~TALENT i~udostępnionej przez TabArena, co wymaga odróżnienia algorytmu publikacji od konkretnej ścieżki integracji.

\subsection{TabM i poziomy łączenia predykcji}
\label{sec:tabm-zespoly}

TabM wykorzystuje sieć gęstą do wytwarzania wielu predykcji, przy czym podmodele współdzielą znaczną część parametrów. Rozwiązanie nawiązuje do BatchEnsemble i~ogranicza konieczność przechowywania całkowicie niezależnej kopii każdej sieci. W~klasyfikacji predykcje podmodeli są łączone w~ocenę klas~\cite{gorishniy2025tabm}.

Oszczędność parametrów nie oznacza zerowego kosztu dodatkowych predykcji. Liczba podmodeli, sposób ich uczenia i~agregacji należą do definicji wykonanego wariantu. W~tej pracy TabM reprezentuje efektywne zespoły neuronowe. Samo zestawienie go z~lokalnym MLP nie jest kontrolą, w~której jedyną różnicą jest uśrednianie: zmienia się również parametryzacja i~procedura uczenia.

Należy odróżnić trzy poziomy łączenia, zestawione w~tabeli~\ref{tab:poziomy-zespolow}. Wewnętrzny zespół TabM występuje wewnątrz pojedynczego dopasowania modelu. Uśrednienie foldów łączy modele uczone na różnych częściach danych. Zespół konfiguracji łączy z~kolei wyniki różnych kandydatów, z~wagami ustalonymi na walidacji. Poziomy mogą się nakładać; nie są wzajemnie wykluczającymi się warunkami.

\begin{table}[tbp]
\centering
\caption{Poziomy łączenia predykcji. Opis definicji metod i przyjętego protokołu, bez własnych wyników jakości.}
\label{tab:poziomy-zespolow}
\small
\begin{tabularx}{\textwidth}{@{}P{0.22\textwidth}YY@{}}
\toprule
Poziom & Łączone obiekty & Znaczenie dla badania \\
\midrule
Wewnętrzny & Podmodele ze współdzieleniem parametrów & Składnik architektury TabM~\cite{gorishniy2025tabm}. \\
Foldy & Modele tej samej konfiguracji & Zwykłe warunki D i T już wykorzystują takie uśrednienie. \\
Konfiguracje & Wyniki różnych kandydatów & T+E dodaje dobór wag na OOF, bez selekcji na teście. \\
\bottomrule
\end{tabularx}
\end{table}

Dla metod strojonych z~ośmiofoldową walidacją konfiguracja domyślna (D) i~wybrana konfiguracja (T) korzystają z~uśrednienia predykcji ośmiu foldów. Warunek T+E dodatkowo łączy konfiguracje. Nie są to trzy niezależne kampanie treningowe ani porównanie jednej sieci z~ośmioma. Natywny IID KNN i~stały TabPFN mają odmienne zasady, które należy zachować w~opisie wyników. Rysunek~\ref{fig:zespoly} ma uwidocznić właśnie to rozróżnienie.

\plannedfigure
    {figures/ch03/poziomy-zespolow.pdf}
    {Wewnętrzny zespół, uśrednienie foldów i łączenie konfiguracji. Opracowanie własne na podstawie mechanizmu TabM~\cite{gorishniy2025tabm} i przyjętego protokołu.}
    {fig:zespoly}
    {Trzy panele z podpisanymi jednostkami łączenia. Zaznaczyć, że D i T zwykłych metod już uśredniają foldy, a KNN i TabPFN mają wyjątki.}

\subsection{TabPFN-3.5: pretrenowanie i kontekst zadania}
\label{sec:tabpfn}

TabPFN-3.5 wykorzystuje wagi uzyskane przed otrzymaniem konkretnego zadania. Oznakowane obserwacje uczące stanowią kontekst, na podstawie którego model przewiduje etykiety nowych przypadków. Lokalna adaptacja nie jest utożsamiana z~trenowaniem całej sieci od początku. Raport techniczny~\cite{jaeger2026tabpfn35} rozróżnia model bazowy oraz warianty o~innych możliwościach i~kosztach inferencji.

W~pracy przyjęto jeden lokalny checkpoint bazowy, ze stałymi ustawieniami inferencji, bez HPO i~bez ośmiofoldowego zespołu treningowego. Kontekst pochodzi z~pełnej części deweloperskiej. Wewnętrzne ustawienie liczby estymatorów jest elementem tej stałej procedury i~nie tworzy warunku T+E. Warianty Plus lub Thinking nie są zamiennikami wybranego checkpointu.

Porównanie kosztu dotyczyć będzie lokalnego przygotowania i~predykcji, bez historycznego kosztu pretrenowania. Wcześniejsze publikacje TabPFN nie potwierdzają automatycznie możliwości ani wyników wersji 3.5. Ponadto zbiorcze zestawienia rodziny w~raporcie technicznym nie muszą dotyczyć jednego wariantu bazowego~\cite{jaeger2026tabpfn35}. Faktyczna dostępność checkpointu, obsługa rozmiarów i~klas oraz wymagania pamięciowe wymagają kwalifikacji. Ich opis nie będzie zastąpiony wynikami autorów raportu lub testem samego interfejsu.

\subsection{Granice porównań i uzasadnienie wyłączeń}
\label{sec:granice-metod}

Przyjęta pula obejmuje odmienne mechanizmy, ale nie całość literatury. ResNet, NODE, podejścia samonadzorowane, inne modele uwagi i~kolejne modele fundamentowe pozostają kontekstem dla wyboru. ResNet jest ważnym odniesieniem historycznym, lecz rozwinięty model gęsty reprezentuje w~macierzy RealMLP. ModernNCA reprezentuje uczenie sąsiedztwa. FT-Transformer wykorzystuje uwagę nad tokenami cech, a~TabM --- współdzielenie parametrów zespołu. Dodawanie kolejnych podobnych ról zwiększałoby koszt bez automatycznego pogłębienia odpowiedzi na przyjęte pytania.

Tak ograniczona pula pozwala zbadać wybrane sposoby wykorzystania sieci przy dostępnym budżecie obliczeń. Wnioski z~porównania będą dotyczyły tych konkretnych procedur i~warunków oceny. Wpływ ważenia klas i~regularyzacji zostanie oceniony osobno, zgodnie z~projektem przedstawionym w~kolejnym rozdziale.
\clearpage
\section{Metodyka badań}
\label{chap:metodyka}

Metodyka obejmuje porównanie kompletnych procedur predykcyjnych oraz dwa badania kontrolowane: ważenia klas i~regularyzacji MLP. Wspólny podział zewnętrzny wyznacza populację oceny, natomiast walidacja wewnętrzna służy selekcji konfiguracji i~łączeniu predykcji.

\subsection{Projekt badania i jednostki analizy}
\label{sec:projekt-badania}

Badanie główne obejmuje siedem procedur neuronowych oraz cztery odniesienia konwencjonalne opisane w~rozdziale~\ref{chap:metody}. Każda para zbiór--metoda podlega ocenie na wspólnym zewnętrznym zbiorze testowym danego zadania. Porównaniu podlegają kompletne procedury, obejmujące przygotowanie cech, uczenie, wybór hiperparametrów i~łączenie predykcji. Różnica między dwiema takimi procedurami nie identyfikuje wpływu pojedynczego składnika architektury.

Jednostką porównania między zadaniami jest zbiór danych. Foldy wewnętrzne, konfiguracje hiperparametrów i~warunki raportowania są zależnymi elementami procedury ocenianej na tym samym zadaniu; nie zwiększają liczby niezależnych jednostek~\cite{demsar2006comparisons}. W~badaniach mechanizmów jednostką dopasowania porównań jest natomiast para warunków o~wspólnym podziale, podzbiorze i~ziarnie. Rozdzielenie tych poziomów określa zakres agregacji i~interpretację niepewności w~sekcji~\ref{sec:analiza-plan}.

Dwa badania uzupełniające wprowadzają kontrolowane zmiany przy zachowaniu pozostałych ustawień. Badanie~A dotyczy ważenia klas, a~badanie~B --- regularyzacji MLP i~liczby przykładów treningowych. Interpretacja tych porównań pozostaje odrębna od rankingu badania głównego. Ich wyniki testowe nie stanowią podstawy zmiany porównywanych procedur.

\subsection{Dobór i audyt zbiorów danych}
\label{sec:dane-audyt}

Dobór dwunastu zadań jest celowy. Ma uwzględnić różne liczności, typy cech, liczbę klas oraz relację liczby cech do liczby przykładów. Tabela~\ref{tab:rejestr-zadan} przedstawia deklarowane źródła i~role danych. Wskazane własności uzasadniają dobór, lecz nie są zmierzoną oceną trudności ani dowodem niezależności źródeł. HIGGS oznacza wskazany podzbiór OpenML, a~nie pełny pierwotny zbiór. Wersja~1 przypisana konfiguracji źródła UCI jest oznaczeniem lokalnej definicji; nie gwarantuje niezmienności pobieranego archiwum.

\begin{table}[tbp]
\centering\small
\caption{Źródła i~role przyjętych zadań. Liczba klas wynika z~definicji zadania.}
\label{tab:rejestr-zadan}
\begin{tabularx}{\linewidth}{P{0.23\linewidth}P{0.23\linewidth}cY}
\toprule
Zbiór & Źródło i~wersja & Klasy & Rola w~projekcie \\
\midrule
Adult & OpenML 1590, v2 & 2 & Cechy mieszane \\
Covertype & OpenML 1596, v4 & 7 & Duża liczność, wiele klas \\
Credit Card Fraud & OpenML 1597, v1 & 2 & Rzadkie zdarzenie \\
HIGGS, podzbiór & OpenML 23512, v2 & 2 & Cechy numeryczne \\
Amazon Employee Access & OpenML 4135, v1 & 2 & Kategorie, rzadka decyzja \\
Santander Customer Satisfaction & Oficjalny lokalny plik treningowy & 2 & Niezrównoważenie klas \\
Online Shoppers & UCI 468 & 2 & Cechy mieszane \\
Diabetes 130-US & UCI 296 & 3 & Powtarzające się pobyty pacjentów \\
Helena & OpenML 41169, v1 & 100 & Wiele klas, ograniczona semantyka cech \\
German Credit & UCI 144 & 2 & Mała próba, cechy mieszane \\
Ionosphere & UCI 52 & 2 & Mała próba, cechy numeryczne \\
Bioresponse & OpenML 4134, v1 & 2 & Dużo cech względem liczby przykładów \\
\bottomrule
\end{tabularx}
\end{table}

Definicja zadania obejmuje znaczenie etykiety, populację predykcji, mapowanie klas, typy kolumn i~sposób oznaczania braków. Ocena poprawności danych uwzględnia stałe cechy, duplikaty, identyfikatory oraz dostępność predyktorów w~momencie podejmowania decyzji. Reguły wyłączeń są niezależne od wyników porównywanych modeli; zmiana populacji na podstawie jakości testowej wprowadzałaby dodatkowy poziom selekcji.

W~Diabetes identyfikator pacjenta służy grupowaniu, a~identyfikatory pacjenta i~pobytu nie są predyktorami. W~Santander identyfikator wiersza jest odrębny od cech. Dla Credit Card Fraud istotne są interpretacja kolumny \texttt{Time} i~ograniczenia losowego podziału, dla Amazon --- powtarzające się encje, a~dla Covertype --- możliwe zależności przestrzenne. Anonimowość cech Heleny ogranicza ocenę ich znaczenia i~dostępności. Sama zgodność schematu pliku nie rozstrzyga tych kwestii.

\subsection{Podziały i przepływ informacji}
\label{sec:podzialy}

W~badaniu głównym wykorzystuje się pierwszy podział stratyfikowanej procedury trójfoldowej z~tasowaniem i~ziarnem~42. Około dwie trzecie obserwacji przypada na część przeznaczoną na uczenie i~walidację, dalej nazywaną częścią rozwojową, a~jedna trzecia na test. Pozostałe dwa podziały zewnętrzne nie są wykorzystywane jako powtórzenia eksperymentu. W~części rozwojowej zwykłe procedury uczone wykorzystują osiem foldów wewnętrznych w~jednym powtórzeniu. Automatyczne przechodzenie do powtarzanej walidacji dla małych danych zostało wyłączone.

Predykcje OOF (ang. out-of-fold) są wyznaczane dla obserwacji spoza części treningowej danego foldu. W~lokalnych MLP i~FT-Transformer etykiety walidacyjne nie uczestniczą w~obliczaniu gradientów, ale służą wyborowi epoki i~wcześniejszemu zatrzymaniu uczenia. OOF nie jest więc oceną na danych całkowicie niewykorzystanych przy dopasowaniu procedury. Predykcje te służą do wyboru konfiguracji i~wag zespołu. Zewnętrzny zbiór testowy służy niezależnej ocenie wybranej procedury~\cite{cawley2010selection}.

Dla Diabetes stosuje się stratyfikowane podziały grupowe na obu poziomach. Wszystkie pobyty jednego pacjenta muszą należeć do tej samej strony danego podziału. Kryteria poprawności podziału obejmują rozłączność grup, obecność klas oraz dokładnie jednokrotne przypisanie obserwacji do walidacji wewnętrznej. Udziały i~proporcje klas przy podziale grupowym mogą różnić się od idealnie równych; raport obejmuje rzeczywiste liczności.

Schemat~\ref{fig:podzial-selekcja} przedstawia dopuszczalny przepływ informacji. Jednokrotny podział ogranicza koszt, ale nie pozwala oszacować zmienności wynikającej ze zmiany podziału. Brak procedury czasowej oznacza również, że wyniki nie stanowią dowodu odporności na zmianę rozkładu w~czasie.

\plannedfigure
{figures/ch04/podzial-i-selekcja.pdf}
{Schemat podziału danych, wyboru konfiguracji i~końcowej oceny; opracowanie własne.}
{fig:podzial-selekcja}
{\textbf{Część rozwojowa}\quad (około $2/3$ obserwacji)\\
Osiem foldów $\longrightarrow$ predykcje OOF\\
$\longrightarrow$ wybór konfiguracji i~wag zespołu\\[0.8em]
\textbf{Część testowa}\quad (około $1/3$ obserwacji)\\
Predykcje wybranej procedury $\longrightarrow$ ocena jakości\\[0.8em]
Etykiety testowe nie uczestniczą w~selekcji.\\
IID KNN: LOO.\quad Stały TabPFN: pełny kontekst rozwojowy, bez OOF.}

\subsection{Przygotowanie cech i wyjątki proceduralne}
\label{sec:przygotowanie-procedura}

Wspólne dane wejściowe zachowują wartości i~typy cech oraz ustalone mapowanie klas. Nie narzuca się wszystkim metodom identycznego kodowania ani skalowania. Przygotowanie danych stanowi część badanego potoku, a~wyuczone przekształcenia zwykłych modeli są dopasowywane wyłącznie na treningu odpowiedniego foldu. Obsługa braków i~nieznanych kategorii musi być określona przed transformacją walidacji i~testu. Odmienne przetwarzanie kategorii jest właściwością porównywanych procedur, a~nie wyizolowanym czynnikiem eksperymentalnym.

Wyjątek stanowi natywny KNN w~zadaniach traktowanych jako IID. Korzysta on z~predykcji leave-one-out, dalej LOO, pojedynczego modelu potomnego. Etykieta przewidywanej obserwacji zostaje wyłączona z~jej sąsiedztwa, ale przygotowanie cech obejmuje całą rozwojową populację wejściową. Nie jest to ponowne dopasowanie wszystkich przekształceń po usunięciu każdego wiersza. Kategorie KNN kodowane są ustalonym kodowaniem one-hot; efektywna liczba sąsiadów i~ewentualne ograniczenie banku przez budżet należą do metadanych wykonania. W~Diabetes KNN korzysta z~ośmiu foldów grupowych.

Stały TabPFN-3.5 wykorzystuje pełną część rozwojową jako oznakowany kontekst. Nie tworzy OOF, nie przechodzi strojenia i~nie otrzymuje zespołu foldów. Jego osiem estymatorów inferencji nie jest ośmioma modelami walidacji. Porównanie obejmuje lokalny model o~ustalonej wersji wag; wiedza uzyskana podczas pretrenowania pozostaje zewnętrznym zasobem tej procedury.

\subsection{Strojenie, losowość i budżety}
\label{sec:strojenie}

Dla każdej z~dziesięciu metod podlegających strojeniu przyjęto konfigurację domyślną i~25 losowych kandydatów. Przestrzenie sześciu lokalnych integracji pochodzą z~plików YAML; KNN, RealMLP, ModernNCA i~TabM korzystają z~generatorów przypiętej biblioteki. Zapisywany jest efektywny zestaw parametrów po zastosowaniu generatorów i~nadpisań. Rozdzielenie pochodzenia przestrzeni jest istotne dla odtworzenia rzeczywiście ocenionych kandydatów.

Ziarno wyszukiwania wynosi~42. Ziarno zewnętrzne wpływa także na inicjalizację modeli, natomiast ziarna wewnętrznych podziałów są określane przez bibliotekę i~zapisywane w~artefaktach. Procedura nie korzysta z~jednego wspólnego ziarna dla wszystkich operacji. TabPFN ma stałe ziarno inferencji~0. Ziarna badań mechanizmów stanowią odrębny element definicji porównania sparowanego.

Budżet pojedynczej konfiguracji wynosi 3600~s, łącznie dla jej foldów i~przygotowania danych. Jest to limit przekazywany bibliotece, a~nie twardy termin zakończenia procesu. W~sześciu uczonych metodach neuronowych ograniczenie czasu stanowi dodatkowe kryterium zakończenia uczenia, obok maksymalnej liczby epok i~reguł opartych na walidacji. Procedura zatrzymuje uczenie z~zachowaniem stanu modelu dostępnego do predykcji; samo zatrzymanie budżetowe nie oznacza niepowodzenia dopasowania. Kontrola czasu odbywa się wewnątrz procedur uczenia, dlatego rzeczywisty czas może przekroczyć przydział. Porównanie dotyczy jakości osiąganej przy wspólnym nominalnym budżecie, a~nie po jednakowej liczbie aktualizacji parametrów. Jednakowa liczba prób nie zapewnia jednakowego czasu ani równoważnej pojemności przestrzeni.

\subsection{Warunki raportowania D, T i T+E}
\label{sec:warunki-raportowania}

Dla zwykłej konfiguracji predykcję testową stanowi średnia prawdopodobieństw z~ośmiu modeli foldów. Nie następuje dodatkowe dopasowanie wybranej konfiguracji na całej części rozwojowej. Ten sposób wykorzystania zapisanych predykcji jest inspirowany organizacją ocen w~TabArena~\cite{erickson2025tabarena}; niniejsze badanie ma jednak własny zakres i~jawne wyjątki.

\begin{table}[tbp]
\centering\small
\caption{Warunki raportowania oraz konsekwencje wyjątków. Wszystkie decyzje o~wyborze wykorzystują część rozwojową.}
\label{tab:warunki}
\begin{tabularx}{\linewidth}{P{0.14\linewidth}YY}
\toprule
Warunek & Wybór & Predykcja testowa \\
\midrule
D & Konfiguracja domyślna & Średnia foldów; dla IID KNN predykcja natywnego modelu \\
T & Najlepsza jakość OOF spośród udanych kandydatów, włącznie z~domyślnym & Zapisana predykcja wybranej konfiguracji \\
T+E & Zachłanny dobór wag na OOF, budżet do~100 kroków, ziarno~42 & Ważona suma zapisanych prawdopodobieństw konfiguracji \\
Stały TabPFN & Brak wyboru konfiguracji i~wag & Predykcja z~pełnym kontekstem rozwojowym; wyłącznie D \\
\bottomrule
\end{tabularx}
\end{table}

Dla konfiguracji $c$ podlegającej ośmiofoldowej walidacji niech $p_{c,k}(x)$ oznacza predykcję modelu foldu~$k$. Warunki D i~T używają $\bar p_c(x)=\frac{1}{8}\sum_{k=1}^{8}p_{c,k}(x)$, odpowiednio dla konfiguracji domyślnej i~wybranej. Warunek T+E ma postać
\begin{equation}
p_{\mathrm{T+E}}(x)=\sum_{c\in\mathcal{C}}\omega_c\bar p_c(x),
\qquad \omega_c\geq0,\quad \sum_{c\in\mathcal{C}}\omega_c=1,
\label{eq:zespol-konfiguracji}
\end{equation}
gdzie $\mathcal{C}$ jest zbiorem dopuszczonych konfiguracji. Wagi $\omega_c$ są wybierane na OOF. Łączenie konfiguracji dodaje więc poziom agregacji do już istniejącego uśredniania foldów; zespół wewnętrzny TabM pozostaje częścią $p_{c,k}$. Dla natywnego IID KNN w~miejsce $\bar p_c$ używana jest predykcja pojedynczego dopasowania, a~stały TabPFN występuje tylko w~D.

W~przypiętej implementacji zachowywany jest najkrótszy prefiks trajektorii zachłannego wyboru osiągający najlepszą ocenę OOF. Budżet 100 kroków nie oznacza 100 zachowanych wyborów ani 100 różnych konfiguracji; powtórny wybór tego samego kandydata zwiększa jego udział. Różnica T+E względem T mierzy korzyść z~tej dodatkowej procedury selekcji i~agregacji, przy wspólnym koszcie znalezienia kandydatów.

Wybór T i~wag T+E następuje przed wykorzystaniem etykiet testowych do obliczenia metryk. Wynik testowy nie służy wyborowi warunku raportowania. Brak kompletu wymaganych dopasowań i~predykcji lub awaria konfiguracji domyślnej uniemożliwiają pełny raport. Konfiguracja, której uczenie zakończyło się poprawnie wskutek ograniczenia czasu, pozostaje dopuszczona do raportowania oraz wyboru T i~T+E, o~ile spełnia te same wymagania kompletności co pozostali kandydaci. Jawnie nieudane konfiguracje losowe mogą zostać pominięte w~wyborze dopiero po zakończeniu prób; ich liczba i~koszt pozostają częścią opisu kompletności.

\subsection{Metryki i reguły decyzji}
\label{sec:metryki}

Dla dziewięciu zadań binarnych główną metryką wyboru jest ROC AUC, dla której większa wartość oznacza lepsze uporządkowanie. Dla trzech zadań wieloklasowych jest to strata logarytmiczna, dla której mniejsza wartość jest korzystniejsza:
\begin{equation}
L_{\log}=-\frac{1}{n}\sum_{i=1}^{n}\log p_{i,y_i}.
\label{eq:logloss-badanie}
\end{equation}
Wartość dotyczy prawdopodobieństwa przypisanego właściwej klasie. Poprawne obliczenie miary wymaga zgodności kolejności kolumn predykcji z~mapowaniem klas; sama zgodność wymiarów nie wystarcza.

Raport uzupełniają accuracy, balanced accuracy, F1 w~wariancie makro i~ważonym, strata logarytmiczna, metryki poszczególnych klas oraz wieloklasowe AUC one-vs-rest makro tam, gdzie jest określone. Główne etykiety decyzyjne powstają przez \(\arg\max\) prawdopodobieństw. Nie oznacza to optymalizacji kosztu aplikacyjnego. Moduł metryk zapisuje nieokreślone ROC AUC jako brak wartości, lecz nie zapisuje osobnego pola z~przyczyną. Dla precyzji, czułości i~F1 stosuje natomiast konwencję \texttt{zero\_division=0}: przy zerowym mianowniku zwraca zero. Rozróżnienie tych przypadków od zerowej wartości poprawnie określonej miary wymaga uwzględnienia liczności klas i~macierzy pomyłek.

W~badaniu rzadkich klas podstawową miarą jest average precision, dalej AP. W~standardowej definicji sumuje ona przyrosty czułości ważone odpowiadającą im precyzją; nie jest to trapezowe pole pod krzywą precision--recall. Uzupełnienie AUC o~tę perspektywę jest uzasadnione odmiennym zachowaniem ocen przy niezrównoważeniu~\cite{saito2015pr}. Interpretacja AP uwzględnia częstość klasy rzadkiej i~liczbę jej przykładów w~teście.

Rzadkie zdarzenie jest określone przez surową etykietę i~jej zapisane kodowanie. Zbiorcze binarne AP raportera dotyczy klasy zakodowanej jako~1, która nie musi oznaczać zdarzenia rzadkiego. W~badaniu~A kolumna prawdopodobieństw i~metryka odpowiadają semantycznie wskazanej klasie rzadkiej. Techniczne nazwy pól raportu i~ich mapowanie na populacje oceny opisano w~rozdziale~\ref{chap:srodowisko}.

\subsection{Badanie A: ważenie klas}
\label{sec:projekt-wazenie}

Badanie obejmuje Credit Card Fraud, Santander i~Amazon oraz metody MLP, TabM i~CatBoost. Dla każdego zestawu porównuje się uczenie nieważone i~ważone. Wagi klas są wyznaczane wyłącznie z~treningowej części danego foldu:
\begin{equation}
w_c=\frac{N_{\mathrm{tr}}}{K n_{c,\mathrm{tr}}},
\label{eq:wagi-klas}
\end{equation}
gdzie \(K\) oznacza liczbę klas, a~\(n_{c,\mathrm{tr}}\) --- liczbę treningowych przykładów klasy~\(c\). Zmianie podlega strata treningowa. Walidacja i~test zachowują oryginalny rozkład oraz nieważone metryki. Badanie nie obejmuje resamplingu ani odrębnego przeszukiwania hiperparametrów dla obu warunków.

Pozostałe hiperparametry są określane na podstawie wyboru OOF w~badaniu głównym. Porównanie wykorzystuje te same foldy i~trzy sparowane ziarna uczenia. W~ten sposób porównanie odpowiada na pytanie o~interwencję przy ustalonej procedurze, a~nie o~najlepszy wynik dwóch osobno strojonych rodzin. Układ obejmuje 54 oceny warunków i~432 dopasowania foldów: 288 neuronowych oraz 144 CatBoost na CPU.

Zgodnie z~analizą w~sekcji~\ref{sec:trudne-rezimy} wzrost czułości po ważeniu może wynikać z~przesunięcia punktu pracy bez poprawy rankingu. AP, ROC AUC i~nieważona strata logarytmiczna pozwalają rozdzielić te efekty. Precyzja, czułość, liczby FP i~FN oraz udział alarmów podlegają ocenie przy progach określonych na OOF. Wybór progu na walidacji definiuje badawczy punkt pracy; bez kosztów błędów nie stanowi optimum wdrożeniowego. Próg wybrany na OOF jest stosowany do średniej predykcji foldów, dlatego jego przenoszenie między tymi rozkładami ocen stanowi dodatkowe ograniczenie interpretacji.

\subsection{Badanie B: regularyzacja i liczność treningu}
\label{sec:projekt-regularyzacja}

Badanie obejmuje German Credit, Ionosphere i~Bioresponse. Punktem odniesienia jest procedura uczenia MLP oparta na opublikowanym rozwiązaniu i~ustalona niezależnie od wyników strojenia na tych trzech zadaniach. Procedura musi zawierać niezerowy dropout oraz niezerowy zanik wag. Nie należy utożsamiać jej automatycznie z~lokalnym bazowym MLP badania głównego. Przedmiotem porównania są pełna procedura, wariant bez dropoutu i~wariant bez zaniku wag. Brak wariantu usuwającego oba składniki ogranicza wnioski o~ich interakcji.

Dla każdego treningowego foldu tworzy się zagnieżdżone, stratyfikowane podzbiory o~udziałach 25\%, 50\% i~100\%. Redukcja obejmuje wyłącznie trening, z~ponownym dopasowaniem jego przekształceń. Walidacja i~test pozostają stałe. Definicja podzbioru obejmuje rzeczywiste liczności klas po zaokrągleniach; nominalny udział nie wystarcza do odtworzenia małych podzbiorów. Te same indeksy i~trzy sparowane ziarna są stosowane dla wszystkich trzech wariantów regularyzacji.

Projekt~\ref{fig:podzbiory} pozwala oddzielić interwencję regularyzacyjną od zmiany składu walidacji. Nie jest krzywą kosztu etykietowania: nawet najmniejszy trening korzysta z~niezmniejszonej walidacji. Układ obejmuje 81 ocen warunków i~648 neuronowych dopasowań foldów. Nie przeprowadza się odrębnego przeszukiwania 26 konfiguracji dla każdej komórki projektu.

\plannedfigure
{figures/ch04/zagniezdzone-podzbiory.pdf}
{Układ zagnieżdżonych podzbiorów treningowych w~badaniu regularyzacji; opracowanie własne.}
{fig:podzbiory}
{\textbf{Trening wewnętrznego foldu}\\
$\mathcal{T}_{25\%}\subset\mathcal{T}_{50\%}\subset\mathcal{T}_{100\%}$\\[0.8em]
Każdy podzbiór: pełna procedura MLP, bez dropoutu, bez zaniku wag.\\
Te same indeksy i~sparowane ziarna dla wariantów regularyzacji.\\[0.8em]
\textbf{Walidacja i~test}: stałe, poza redukcją treningu.\\
Przekształcenia cech dopasowywane na właściwym podzbiorze treningowym.}

Analiza generalizacji łączy testowe AUC z~historią uczenia oraz nieważoną stratą logarytmiczną na treningu i~walidacji. Luka jest definiowana dla tej samej sieci foldu, w~trybie ewaluacji z~wyłączonym dropoutem. Różnica między stratą treningową pojedynczej sieci a~testową stratą zespołu ośmiu sieci nie jest luką generalizacji tej samej procedury. Mniejsza luka sama w~sobie nie oznacza lepszej generalizacji; może towarzyszyć pogorszeniu obu ocen.

\subsection{Analiza porównawcza i niepewność}
\label{sec:analiza-plan}

Analiza rozdziela wielkość różnicy, jej zmienność oraz zakres wnioskowania. Podstawowy panel obejmuje T dla dziesięciu metod strojonych i~odrębnie oznaczony stały TabPFN. D i~T+E służą do analizy zmiany procedury. Każdy wynik pozostaje przypisany do jednego zadania, metryki i~warunku.

\begin{table}[tbp]
\centering\small
\caption{Powiązanie pytań badawczych z~porównaniami i~jednostkami analizy.}
\label{tab:pytania-analiza}
\begin{tabularx}{\linewidth}{P{0.09\linewidth}YY}
\toprule
Pytanie & Porównanie i~miara & Jednostka i~granica wniosku \\
\midrule
PB1 & Wyniki wszystkich metod; kontrasty między procedurami & Zbiór; celowo wybrane zadania \\
PB2 & D--T i~T--T+E; zmiana metryki oraz koszt & Ta sama para zbiór--metoda \\
PB3 & Ważenie względem kontroli; AP i~decyzje progowe & Sparowane ziarna w~trzech zadaniach \\
PB4 & Ablacja względem pełnej procedury; AUC, strata i~luka & Wspólny podzbiór i~ziarno w~trzech zadaniach \\
PB5 & Jakość oraz rozdzielone składniki kosztu & Zadanie, procedura i~określony sprzęt \\
\bottomrule
\end{tabularx}
\end{table}

Kontrasty o~bezpośredniej interpretacji metodologicznej to RealMLP\allowbreak{}--\allowbreak{}CatBoost, TabM\allowbreak{}--\allowbreak{}CatBoost oraz ModernNCA\allowbreak{}--\allowbreak{}KNN. Dla pary \(A,B\) różnica dodatnia oznacza korzyść~\(A\): \(\Delta_{\mathrm{AUC}}=\mathrm{AUC}_A-\mathrm{AUC}_B\), natomiast \(\Delta_{\log}=L_{\log,B}-L_{\log,A}\). Surowe różnice AUC i~straty nie podlegają uśrednianiu na wspólnej skali. Pomocnicze rangi można obliczać jedynie przy podaniu wspólnego zbioru kompletnych wyników i~reguły remisów.

Wnioskowanie o~przewadze między zadaniami wymaga określenia modelu losowego, rodziny porównań i~sposobu postępowania z~remisami. Test znaków dotyczy częstości kierunku różnicy, nie jej wielkości; równoczesne testowanie wielu par wymaga kontroli wielokrotności. Dobór celowy nie gwarantuje niezależności zadań ani modelu losowego uzasadniającego test. Jeżeli założeń nie da się obronić, analiza ma charakter opisowy. Niewielka liczba zadań ogranicza moc, a~brak istotności nie potwierdza równoważności.

Niepewność różnicy na pojedynczym zadaniu można analizować przez sparowany bootstrap zapisanych predykcji. Przy założeniu IID jednostką losowania jest wiersz, z~warstwowaniem po klasie. Dla Diabetes jednostkę stanowi pacjent ze wszystkimi jego pobytami. Te same repliki muszą służyć obu modelom. Liczba replik z~nieokreśloną metryką stanowi część opisu resamplingu. Takie przedziały są warunkowe względem ustalonego podziału i~wytrenowanych modeli; nie mierzą zmienności ponownego uczenia.

W~badaniach~A i~B przedmiotem analizy są wyniki wszystkich trzech sparowanych ziaren, średnia różnica i~jej zakres. Zakres nie jest przedziałem ufności. Foldy nie są traktowane jako niezależne replikacje, a~trzy zadania każdego badania nie uzasadniają przewagi populacyjnej. Pełne zestawienie jedenastu metod nie wymaga testowania wszystkich 55 par; niezgodne skale metryk pozostają rozdzielone.

Braki i~awarie wymagają jawnych mianowników. Analiza kompletnych par może pomijać najtrudniejsze zadania i~przez to zmieniać obraz porównania. Nieudanym uruchomieniom nie przypisuje się sztucznych metryk. Wielkość różnicy oraz jej znaczenie praktyczne pozostają odrębne od wartości~\(p\); bez kosztów aplikacyjnych nie ustala się wspólnej granicy użyteczności dla AUC, AP i~straty.

\subsection{Koszt obliczeń i odtwarzalność}
\label{sec:koszt-metodyka}

Nominalna macierz główna obejmuje \(12\cdot10\cdot26+12=3132\) zadań konfiguracji, 132 pary zbiór--metoda i~372 wiersze raportu dla przewidzianych warunków. Sześć strojonych procedur neuronowych wymaga nominalnie 14\,976 dopasowań foldów; stała inferencja TabPFN jest liczona osobno. Badania mechanizmów dodają 936 dopasowań neuronowych i~144 na CPU. Są to nominalne liczby dopasowań wynikające z~macierzy badania; rzeczywisty koszt zależy także od czasu pojedynczego dopasowania, powtórzeń i~niepowodzeń.

Koszt wyszukiwania i~suma czasów dopasowania konfiguracji wybranych do użycia opisują różne etapy procedury. Czas \texttt{search\_fit\_seconds} obejmuje udane dopasowania, natomiast \texttt{search\_wall\_seconds} uwzględnia zapisane próby, także zakończone błędem i~powtórzenia. Przerwanie bez znanego końca próby może uniemożliwić obliczenie łącznego czasu rzeczywistego. Brakująca wartość nie oznacza zerowego kosztu.

\texttt{selected\_inference\_seconds} sumuje zapisane przebiegi predykcji konfiguracji o~dodatniej wadze. Nie jest nowym pomiarem końcowej latencji całego zespołu i~nie obejmuje automatycznie wszystkich kosztów składania wyników. Czas wyboru z~pamięci podręcznej raportowany jest osobno. Koszt znalezienia kandydatów jest wspólny dla D, T i~T+E, dlatego sumowanie tych wierszy prowadziłoby do jego wielokrotnego naliczenia.

Opis warunków pomiaru obejmuje CPU, GPU, RAM, VRAM, liczbę wątków, rozmiar partii predykcji i~granice pomiarów. Ustawienia obejmują cztery wątki CPU. Limit partii predykcji~4096 jest przekazywany lokalnym adapterom i~wykorzystywany przez MLP, FT-Transformer oraz TabNet. Dla MLP i~FT-Transformer obowiązuje dodatkowo ograniczenie efektywnym rozmiarem partii uczenia. Natywne KNN, RealMLP, ModernNCA i~TabM oraz stały TabPFN nie otrzymują tego ustawienia; podział predykcji na partie określają te implementacje. W~przyjętym podziale zasobów sześć uczonych metod neuronowych korzysta z~GPU, a~odniesienia konwencjonalne z~CPU. Porównanie kosztu na takim podziale zasobów nie izoluje efektu architektury. Backend zapisuje pomiary pamięci w~bloku \texttt{memory\_usage} rekordu konfiguracji, a~nie w~kolumnach zbiorczego CSV. Śledzenie obejmuje wyłącznie dopasowanie: RSS procesu na CPU oraz pamięć przydzieloną i~rezerwowaną przez PyTorch na GPU. Interpretacja odczytu GPU zależy od znacznika dostępności śledzenia. Pomiary te nie określają pełnego zużycia VRAM procesu, szczytu całej próby ani pamięci testowej predykcji. Ta ostatnia wymaga osobnego pomiaru, szczególnie dla KNN i~TabPFN. Zakres pomiaru nie obejmuje zewnętrznego pretrenowania TabPFN, kosztu finansowego ani emisji.

W~FT-Transformerze efektywny rozmiar partii jest dodatkowo ograniczany na podstawie liczby cech i~głów uwagi. Dla $p$ predyktorów i~$h$ głów przyjmuje się $4h(p+1)^2$ bajtów na przykład dla pojedynczego tensora wyników uwagi w~formacie float32, z~uwzględnieniem tokenu klasyfikacyjnego. Z~przydziału 4~GiB wyznacza się limit liczby przykładów, zaokrąglany w~dół do potęgi dwójki, z~dolnym ograniczeniem tego limitu do ośmiu. Efektywna partia jest mniejszą z~wartości skonfigurowanej i~tak wyznaczonego limitu. Przy 1776 predyktorach limit wynosi 32 przykłady dla ośmiu głów oraz 64 dla czterech. Ograniczenie obowiązuje także podczas walidacji i~predykcji. Jest to heurystyka pamięciowa dotycząca jednego tensora, nie gwarancja zmieszczenia pełnego przebiegu w~VRAM. Ponieważ zmniejszenie partii zmienia również liczbę aktualizacji parametrów na epokę, efektywny rozmiar należy uwzględniać przy interpretacji zarówno kosztu, jak i~jakości modelu.

Tożsamość danych, zależności i~sprzętu wiąże pomiar z~konkretnym wykonaniem. Mechanizm zapisu tej informacji oraz kontroli zgodności opisano w~rozdziale~\ref{chap:srodowisko}.

\clearpage
\section{Środowisko eksperymentalne i odtwarzalność}
\label{chap:srodowisko}

Badanie główne obejmuje nominalnie 3132 zadania konfiguracji i~14\,976 dopasowań sieci uczonych w~foldach. Przy budżecie godziny na konfigurację nawet pojedyncze niepowodzenie może oznaczać istotny koszt powtórzenia obliczeń. Środowisko przygotowano więc tak, aby ukończone wyniki można było zachować przy wznowieniu kampanii, a~raporty odtworzyć z~zapisanych predykcji. Każdy wynik pozostaje powiązany z~danymi, ustawieniami i~warunkami wykonania.

\subsection{Wymagania wynikające z metodyki}
\label{sec:wymagania-platformy}

Platforma musi zapewniać spójne mapowanie danych i~klas, kontrolę podziałów, odtwarzalne generowanie konfiguracji oraz wybór modeli bez korzystania z~etykiet testowych. Powinna ponadto umożliwiać przerwanie długiej kampanii, rozpoznanie kompletnych wyników i~odtworzenie raportów bez ponownego uczenia. Te wymagania wynikają z~liczby konfiguracji i~z~konieczności przypisania każdej obserwacji naukowej do konkretnego wykonania.

Przyjęto lokalne, sekwencyjne uruchamianie konfiguracji i~foldów w~jednym procesie sterującym. Ustawienia z~plików YAML są odczytywane i~sprawdzane przez centralny moduł konfiguracji. Zachowano obiektowe adaptery modeli oraz wspólny kontrakt predykcji prawdopodobieństw. Umożliwia to współistnienie lokalnych implementacji, natywnych modeli bibliotecznych i~stałego modelu pretrenowanego.

Wymaganiem naukowym jest również jawność niepowodzeń. Awaria nie może zostać przekształcona w~nieoznaczony brak lub pozornie niski koszt. Zapis próby obejmuje stan, czas i~diagnostykę, natomiast reguły raportowania odróżniają jawnie nieudaną próbę od nieukończonego zadania.

\subsection{Architektura i integracje}
\label{sec:architektura-platformy}

Za uruchamianie modeli i~walidację odpowiadają biblioteki TabArena i~AutoGluon w~ustalonych wersjach. Własny kod łączy je z~konfiguracją badania, obsługą danych oraz zapisem wyników. Zakres tej integracji i~odpowiedzialność poszczególnych modułów przedstawia tabela~\ref{tab:warstwy-platformy}.

\begin{table}[tbp]
\centering\small
\caption{Warstwy platformy i~ich odpowiedzialność. Nazwy modułów wskazują źródła implementacji, a~nie niezależne usługi.}
\label{tab:warstwy-platformy}
\begin{tabularx}{\linewidth}{P{0.22\linewidth}P{0.27\linewidth}Y}
\toprule
Warstwa & Główne moduły & Odpowiedzialność \\
\midrule
Konfiguracja & \path{configuration.py}, \path{configs/} & Typy, walidacja i~łączenie ustawień z~plików oraz nadpisań \\
Dane & \path{data/tasks.py} & Przygotowanie, manifesty, mapowanie klas i~zewnętrzne podziały \\
Katalog i~przestrzenie & \path{catalog.py}, \path{specification.py} & Możliwości metod i~losowanie konfiguracji \\
Wykonanie & \path{backend.py}, adaptery modeli & Uczenie foldów, OOF i~prawdopodobieństwa testowe \\
Zasoby & \path{runtime.py} & Walidacja urządzenia, wątki i~kontekst dopasowania \\
Zapis i~wznawianie & \path{runner.py}, \path{artifacts.py} & Tożsamość, blokady, próby, integralność i~wznowienie \\
Raportowanie & \path{reporting.py}, \path{evaluation/metrics.py} & Wybór OOF, zespoły oraz raporty JSON i~CSV \\
\bottomrule
\end{tabularx}
\end{table}

Główna ścieżka wykorzystuje sześć lokalnych integracji: MLP, TabNet, FT-Transformer, regresję logistyczną, XGBoost i~CatBoost. KNN, RealMLP, ModernNCA oraz TabM są uruchamiane przez natywne ścieżki przypiętych bibliotek. Osobny adapter obsługuje stały TabPFN. Wspólny katalog eksperymentów łączy te ścieżki, dzięki czemu dodanie natywnej integracji nie wymaga powielania jej procedury uczenia w~lokalnym adapterze.

Moduł wykonawczy otrzymuje dane z~zachowaniem ról cech numerycznych i~kategorycznych. Automatyczne wspólne przygotowanie cech jest wyłączone, a~za przekształcenia odpowiada właściwa integracja. Przy GPU pomiar dopasowania i~predykcji jest otoczony synchronizacją urządzenia. Ogranicza to błąd wynikający z~asynchronicznego wykonania operacji, ale nie usuwa konieczności precyzyjnego określenia granic pomiaru.

Rysunek~\ref{fig:architektura} przedstawia zależności komponentów i~oddzielenie uczenia od wtórnego odczytu zapisanych predykcji. Interfejs wiersza poleceń znajduje się w~\path{experiments/cli.py}; konfiguracja jest rozwiązywana przed wywołaniem właściwego polecenia.

\begin{figure}[tbp]\centering\small
% Opracowanie własne; schemat procedury opisanej w rozdziałach 3–5.
\begin{minipage}[t]{0.31\linewidth}\centering\textbf{Przygotowanie}\par\medskip
\fbox{\parbox{0.85\linewidth}{\centering Konfiguracja YAML}}\par\smallskip$\downarrow$\par\smallskip\fbox{\parbox{0.85\linewidth}{\centering Dane i~podziały}}\par\smallskip$\downarrow$\par\smallskip\fbox{\parbox{0.85\linewidth}{\centering Katalog konfiguracji}}\end{minipage}
\hfill
\begin{minipage}[t]{0.31\linewidth}\centering\textbf{Wykonanie}\par\medskip
\fbox{\parbox{0.85\linewidth}{\centering TabArena i~AutoGluon}}\par\smallskip$\downarrow$\par\smallskip\fbox{\parbox{0.85\linewidth}{\centering Adapter wybranej metody}}\par\smallskip$\downarrow$\par\smallskip\fbox{\parbox{0.85\linewidth}{\centering Uczenie i~predykcje}}\end{minipage}
\hfill
\begin{minipage}[t]{0.31\linewidth}\centering\textbf{Trwały zapis}\par\medskip
\fbox{\parbox{0.85\linewidth}{\centering Manifest i~historia prób}}\par\smallskip$\downarrow$\par\smallskip\fbox{\parbox{0.85\linewidth}{\centering Predykcje i~metadane}}\par\smallskip$\downarrow$\par\smallskip\fbox{\parbox{0.85\linewidth}{\centering Raport D/T/T+E}}\end{minipage}\par\medskip Przepływ: przygotowanie $\longrightarrow$ wykonanie $\longrightarrow$ zapis.\par Wznowienie odczytuje ukończone zadania; raportowanie korzysta z~zapisanych predykcji.

\caption{Schemat lokalnej platformy oraz przepływu artefaktów; opracowanie własne.}
\label{fig:architektura}
\end{figure}

\subsection{Izolacja informacji i identyfikowalność danych}
\label{sec:izolacja-platformy}

Przygotowane dane są przechowywane jako plik Parquet z~manifestem i~sumami kontrolnymi. Źródła OpenML mają określone identyfikatory i~wersje; tożsamość pobieranych archiwów UCI określa suma kontrolna zawartości. Santander jest pobierany automatycznie z~OpenML (identyfikator 46859, wersja~2) w~ramach wspólnej procedury przygotowania danych. Zgodność schematu i~hashu chroni przed niezamierzoną podmianą danych, lecz nie rozstrzyga problemów populacji, przecieku lub duplikatów wskazanych w~sekcji~\ref{sec:dane-audyt}.

Zewnętrzne podziały zachowują identyfikatory wierszy, etykiety i~grupy pacjentów. Kolumny grupujące nie wchodzą do predyktorów. Backend zapisuje rzeczywiste indeksy walidacji wewnętrznej, co umożliwia późniejszą kontrolę pokrycia OOF i~rozłączności pacjentów. Zdefiniowane są także zastrzeżone nazwy pomocniczych kolumn; konflikt z~danymi wejściowymi jest odrzucany zamiast prowadzić do cichego nadpisania cechy.

Łączenie predykcji wymaga zgodności zarówno obserwacji, jak i~znaczenia kolumn. Macierze o~tych samych wymiarach mogą bowiem opisywać inne wiersze lub mieć odwróconą kolejność klas. Mapowanie klas pochodzi z~definicji przygotowanego zadania i~jest zapisywane wraz z~wynikiem. Podczas zapisu sprawdza się zgodność etykiet z~identyfikatorami wierszy, wymiary macierzy oraz skończoność, zakres i~sumy prawdopodobieństw. Niewielkie odchylenia numeryczne mogą być normalizowane. Przed zbudowaniem zespołu porównuje się ponadto etykiety i~identyfikatory wierszy OOF oraz testu. Raporter nie porównuje osobno list klas między rekordami, dlatego poprawność ich kolejności opiera się na wspólnym mapowaniu zadania; sama kontrola numeryczna jej nie potwierdza.

Ustawienie pola czasowego w~definicji zadania prowadzi do jawnego odrzucenia niewspieranego podziału czasowego. Kontrola techniczna blokuje zatem niewspieraną procedurę, ale nie zastępuje naukowego uzasadnienia podziału IID.

Raport zachowuje odziedziczoną nazwę pola \texttt{log\_loss\_val}, chociaż zawarta w~nim wartość jest stratą na zewnętrznym teście. Przy przenoszeniu wyników do tabel pracy rozstrzygająca jest populacja powstania predykcji. Podobnie zbiorcze binarne AP odnosi się do klasy zakodowanej jako~1; do oceny rzadkiej klasy służy AP zapisane osobno dla każdej klasy w~\texttt{per\_class}. Kolumnę wybiera się na podstawie mapowania rzeczywistej etykiety zdarzenia, a~nie założenia, że ma ona kod~1. Te rozbieżności między nazwą techniczną a~interpretacją pomiaru muszą być uwzględnione w~analizie.

\subsection{Artefakty, awarie i wznowienie}
\label{sec:artefakty-wznowienie}

Artefakty zestawione w~tabeli~\ref{tab:artefakty} pozwalają wznowić obliczenia bez pomieszania wyników różnych konfiguracji. Identyfikator badania jest wyznaczany z~efektywnych ustawień, ścieżek, manifestów danych, skrótu źródeł i~pliku zależności, wybranych wersji pakietów oraz zapisanej charakterystyki sprzętu. Dane są serializowane w~ustalonej postaci i~identyfikowane przez SHA-256. Skrót źródeł obejmuje moduły wykonawcze, \path{pyproject.toml} i~\path{uv.lock}; nie obejmuje automatycznie dokumentacji, testów ani skryptów wdrożeniowych. Identyfikator commitu i~obrazu kontenera należy zachować osobno. Zmiana ścieżek może zmienić tożsamość mimo tej samej treści danych.

\begin{table}[tbp]
\centering\small
\caption{Podstawowe artefakty jednego badania. Ścieżki odnoszą się do katalogu jego tożsamości.}
\label{tab:artefakty}
\begin{tabularx}{\linewidth}{P{0.46\linewidth}Y}
\toprule
Plik lub katalog & Znaczenie \\
\midrule
\path{manifest.json}, \path{index.json} & Zamrożone ustawienia i~wykaz zadań \\
\path{splits/DATASET.json} & Tożsamość zewnętrznego podziału \\
\path{jobs/DATASET-MODEL-000/manifest.json} & Konkretna konfiguracja \\
\path{jobs/.../status.json}, \path{attempts/UUID.json} & Stan i~historia prób \\
\path{jobs/.../records/result.json} & Metadane wykonania oraz czasy \\
\path{jobs/.../predictions/result.npz} & Wyrównane predykcje OOF i~testowe; TabPFN bez OOF \\
\path{failure.json}, \path{failures.json} & Diagnostyka niepowodzeń we właściwym zakresie \\
\path{report.json}, \path{report.csv} & Pochodne wyniki D, T i~T+E \\
\bottomrule
\end{tabularx}
\end{table}

Sama konfiguracja nie wystarcza do odtworzenia warunków uczenia, ponieważ rozmiar partii zależy również od danych danego foldu i~ograniczeń wykonania. Dlatego dla lokalnych MLP, FT-Transformera i~TabNet w~bloku \texttt{hyperparameters\_fit} zapisuje się efektywny rozmiar partii (\texttt{effective\_batch\_size}), a~dla TabNet także skorygowany rozmiar partii wirtualnej (\texttt{effective\_virtual\_batch\_size}). Wartości te pozwalają ustalić, z~jakimi partiami rzeczywiście pracował model po zastosowaniu reguł z~sekcji~\ref{sec:strojenie}.

Ten sam blok zawiera znacznik zakończenia z~powodu budżetu (\texttt{budget\_truncated}), liczbę ukończonych epok (\texttt{completed\_epochs}), przydzielony czas (\texttt{granted\_seconds}) i~czas dopasowania mierzony przez lokalną integrację (\texttt{fit\_seconds\_wrapper}). Zapis obejmuje więc zarówno ustawienia wykonania, jak i~przebieg uczenia. Czas lokalnego dopasowania dotyczy pojedynczego foldu; koszt całej konfiguracji i~wyszukiwania rozlicza się zgodnie z~sekcją~\ref{sec:koszt-metodyka}.

Metadane JSON i~predykcje są najpierw zapisywane do plików tymczasowych. Po opróżnieniu buforów zapisu i~synchronizacji z~dyskiem plik tymczasowy atomowo zastępuje plik docelowy. Blokada wyłączna zapobiega jednoczesnemu zapisowi tego samego katalogu przez dwa uruchomienia. Przed pominięciem ukończonego zadania moduł wykonawczy kontroluje spójność statusu, rekordów, predykcji i~ich sum kontrolnych. Sam fakt istnienia pliku nie wystarcza do uznania zadania za ukończone.

Awarie są zapisywane wraz z~diagnostyką, a~niezależne zadania mogą być kontynuowane. Wystąpienie niepowodzeń powoduje niezerowy kod zakończenia głównego uruchomienia. Automatyczny raport powstaje po jego pomyślnym ukończeniu. Oddzielne polecenie raportowania może obsłużyć jawnie nieudane konfiguracje losowe, gdy wszystkie zadania zostały podjęte; nieudany wariant domyślny albo zadanie nieukończone uniemożliwiają pełny raport.

Po próbie usuwane są dopasowane modele, a~pozostają predykcje i~metadane. Umożliwia to ponowne wyliczenie raportu bez zachowania wszystkich sieci. Nie umożliwia jednak wznowienia w~środku foldu ani wykorzystania usuniętego modelu do nowych predykcji. Konfigurację przerwaną bez kompletnego zapisu wyników trzeba dopasować ponownie w~całości; poprawne zakończenie uczenia wskutek ograniczenia budżetu nie wymaga takiego powtórzenia. Tożsamość badania i~kontrole zgodności ograniczają również automatyczne przenoszenie wyników do zmienionego eksperymentu.

\subsection{Dobór partii i ograniczenia pamięci}
\label{sec:partie-platformy}

W~MLP, FT-Transformerze i~TabNet rozmiar partii wynika z~liczności treningowej części danego foldu. Strojenie obejmuje mnożnik tej bazowej wartości, wybierany ze zbioru $\{0{,}5;\,1;\,2\}$; w~konfiguracji domyślnej wynosi on~1. Dla MLP i~FT-Transformera podstawą doboru jest liczba przykładów treningowych~$n$. TabNet uwzględnia dodatkowo kategorie: progi stosuje się do wartości $n'=n(1+C/5000)$, gdzie $C$ oznacza sumę liczności słowników kategorii utworzonych na treningu danego foldu, wraz z~pozycjami zarezerwowanymi dla nieznanych wartości. Przyjęte progi zestawiono w~tabeli~\ref{tab:dobor-partii}. Jest to heurystyka wykonania, której nie należy utożsamiać z~empirycznie potwierdzonym optymalnym doborem partii.

\begin{table}[tbp]
\centering\small
\caption{Bazowe rozmiary partii w~lokalnych integracjach. Próg odnosi się do $n$ dla MLP i~FT-Transformera oraz do $n'$ dla TabNet.}
\label{tab:dobor-partii}
\begin{tabular}{lrrr}
\toprule
Przedział & MLP i~FT-Transformer & TabNet & Partia wirtualna TabNet \\
\midrule
$[0, 2000)$ & 64 & 128 & 128 \\
$[2000, 8000)$ & 128 & 256 & 128 \\
$[8000, 32000)$ & 256 & 1024 & 128 \\
$[32000, 128000)$ & 512 & 2048 & 256 \\
$[128000, \infty)$ & 1024 & 4096 & 512 \\
\bottomrule
\end{tabular}
\end{table}

Po zastosowaniu mnożnika rozmiar partii zaokrągla się w~dół do potęgi dwójki, z~minimum ośmiu przykładów. Następnie uwzględnia się ograniczenie pamięci FT-Transformera opisane w~sekcji~\ref{sec:koszt-metodyka} lub korektę TabNet zapobiegającą pozostawieniu jednoelementowej partii. W~TabNet osobno dobierana jest partia wirtualna, czyli rozmiar fragmentów partii wykorzystywanych przy normalizacji. Jej wartość bazowa pochodzi z~tej samej tabeli, lecz nie jest skalowana mnożnikiem. Ogranicza się ją do rozmiaru partii uczenia, po czym w~razie potrzeby zwiększa, aby uniknąć jednoelementowych fragmentów także w~ostatniej, niepełnej partii. Jawne zadanie całkowitego rozmiaru partii wyłącza automatyczny dobór i~działanie mnożnika, ale pozostawia wymienione ograniczenia wykonania.

W~FT-Transformerze efektywny rozmiar partii jest dodatkowo ograniczany na podstawie liczby cech i~głów uwagi. Dla $d$ predyktorów i~$H$ głów przyjmuje się $4H(d+1)^2$ bajtów na przykład dla pojedynczego tensora wyników uwagi w~formacie float32, z~uwzględnieniem tokenu klasyfikacyjnego. Z~przydziału 4~GiB wyznacza się limit liczby przykładów, zaokrąglany w~dół do potęgi dwójki, z~dolnym ograniczeniem tego limitu do ośmiu. Efektywna partia jest mniejszą z~wartości uzyskanej po automatycznym doborze i~zastosowaniu mnożnika oraz tak wyznaczonego limitu. Jeżeli rozmiar zadano bezpośrednio jako liczbę całkowitą, limit nakłada się na tę liczbę. Przy 1776 predyktorach limit wynosi 32 przykłady dla ośmiu głów oraz 64 dla czterech. Ograniczenie obowiązuje także podczas walidacji i~predykcji. Jest to heurystyka pamięciowa dotycząca jednego tensora, nie gwarancja zmieszczenia pełnego przebiegu w~VRAM. Ponieważ zmniejszenie partii zmienia również liczbę aktualizacji parametrów na epokę, efektywny rozmiar należy uwzględniać przy interpretacji zarówno kosztu, jak i~jakości modelu.

\subsection{Pomiar kosztu i interpretacja raportu}
\label{sec:pomiary-platformy}

Koszt wyszukiwania i~suma czasów dopasowania konfiguracji wybranych do użycia opisują różne etapy procedury. Czas \texttt{search\_fit\_seconds} obejmuje udane dopasowania, natomiast \texttt{search\_wall\_seconds} uwzględnia zapisane próby, także zakończone błędem i~powtórzenia. Przerwanie bez znanego końca próby może uniemożliwić obliczenie łącznego czasu rzeczywistego. Brakująca wartość nie oznacza zerowego kosztu.

\texttt{selected\_inference\_seconds} sumuje zapisane przebiegi predykcji konfiguracji o~dodatniej wadze. Nie jest nowym pomiarem końcowej latencji całego zespołu i~nie obejmuje automatycznie wszystkich kosztów składania wyników. Czas wyboru z~pamięci podręcznej raportowany jest osobno. Koszt znalezienia kandydatów jest wspólny dla D, T i~T+E, dlatego sumowanie tych wierszy prowadziłoby do jego wielokrotnego naliczenia.

Warunki wykonania obejmują cztery wątki CPU oraz limit partii predykcji 4096, przekazywany lokalnym MLP, FT-Transformerowi i~TabNet. W~dwóch pierwszych metodach obowiązuje również ograniczenie partią uczenia. Natywne integracje i~TabPFN określają podział predykcji na partie we własnych implementacjach. Zakres pomiarów zestawiono w~tabeli~\ref{tab:zakres-pomiarow}.

\begin{table}[tbp]
\centering\small
\caption{Zakres pomiarów pamięci podczas dopasowania.}
\label{tab:zakres-pomiarow}
\begin{tabularx}{\linewidth}{P{0.18\linewidth}YY}
\toprule
Miara & Zakres & Granica interpretacji \\
\midrule
RSS & Pamięć rezydentna procesu na CPU & Nie jest pomiarem całej maszyny \\
Pamięć PyTorch & Pamięć przydzielona i~rezerwowana na GPU & Nie obejmuje wszystkich alokacji procesu ani innych użytkowników GPU \\
Dostępność śledzenia & Znacznik w~bloku \texttt{memory\_usage} & Brak pomiaru nie oznacza zerowego zużycia \\
\bottomrule
\end{tabularx}
\end{table}

Pomiary pamięci są zapisywane w~bloku \texttt{memory\_usage} rekordu konfiguracji i~obejmują dopasowanie. Nie wyznaczają szczytu całej próby ani pamięci testowej predykcji. Osobny pomiar inferencji jest szczególnie istotny dla KNN i~TabPFN, które zachowują zbiór odniesienia lub kontekst zadania. Zewnętrzne pretrenowanie TabPFN pozostaje poza kosztem lokalnego wykonania.

Współdzielenie GPU wymaga dodatkowej kontroli: zapis modelu urządzenia i~jego pamięci nie określa obciążenia przez inne procesy. Przy interpretacji czasu oraz niepowodzeń należy powiązać próby z~dostępnymi zapisami warunków wykonania. Sam komunikat o~wyczerpaniu pamięci nie pozwala przypisać całego zużycia badanemu modelowi.

Moduł metryk zapisuje nieokreślone ROC AUC jako brak wartości, lecz nie zapisuje osobnego pola z~przyczyną. Dla precyzji, czułości i~F1 stosuje natomiast konwencję \texttt{zero\_division=0}: przy zerowym mianowniku zwraca zero. Rozróżnienie tych przypadków od zerowej wartości poprawnie określonej miary wymaga uwzględnienia liczności klas i~macierzy pomyłek.

\subsection{Weryfikacja poprawności platformy}
\label{sec:kwalifikacja}

Weryfikacja platformy dotyczy trzech poziomów: kontraktów danych, poprawności przebiegu eksperymentu i~integralności zapisanych artefaktów. Kontrole wymiarów i~prawdopodobieństw wykrywają błędy interfejsów, natomiast zgodność identyfikatorów wierszy i~grup pozwala sprawdzić przepływ informacji między treningiem, walidacją i~testem.

Test syntetyczny \texttt{smoke} obejmuje niewielki przebieg z~regresją logistyczną i~dwuepokowym MLP. Służy kontroli uruchamiania, zapisu, wznowienia i~raportowania. Jego kryterium stanowi spójność artefaktów i~przebiegu, a~nie jakość predykcji. Test wznowienia wymaga zachowania ukończonych rekordów i~ponowienia wyłącznie tych konfiguracji, dla których nie istnieje kompletny, zgodny zapis.

Poprawność podziału grupowego oznacza rozłączność pacjentów na każdym poziomie, zachowanie wymaganych klas oraz dokładnie jednokrotne pokrycie obserwacji predykcjami OOF. Sam brak wspólnych indeksów wierszy nie wystarcza, ponieważ różne pobyty mogą dotyczyć tej samej osoby. Analogicznie numeryczna poprawność macierzy prawdopodobieństw nie zastępuje zgodności znaczenia jej kolumn.

Dowód weryfikacji odnosi się do konkretnej rewizji programu, danych i~środowiska. Zmiana tych składników ogranicza możliwość przenoszenia wcześniejszych sprawdzeń na nowe wykonanie. Szczególne znaczenie mają granice pamięci, asynchroniczność GPU oraz zależność kosztu inferencji TabPFN od rozmiaru kontekstu.

\subsection{Środowisko wykonania i procedura kampanii}
\label{sec:wykonanie-kampanii}

Podstawą odtworzenia środowiska jest Python~3.11 na Linux x86-64 oraz plik \path{uv.lock}. Istotne deklarowane wersje zestawiono w~tabeli~\ref{tab:wersje}. Wersje pakietów identyfikują implementacje i~nie są tożsame z~datami publikacji metod. Wersja pakietu \texttt{tabpfn} nie zastępuje identyfikacji konkretnego checkpointu modelu.

\begin{table}[tbp]
\centering\small
\caption{Wybrane przypięcia zależności środowiska. Pełna tożsamość wymaga zachowania pliku blokady i~manifestu uruchomienia.}
\label{tab:wersje}
\begin{tabularx}{\linewidth}{YY}
\toprule
Składnik & Wersja \\
\midrule
TabArena & 0.1.0 \\
scikit-learn & 1.9.0 \\
XGBoost & 3.2.0 \\
CatBoost & 1.2.10 \\

AutoGluon core i~tabular & 1.6.3b20260918 \\
PyTorch & 2.13.0 \\
\texttt{pytabkit} & 1.8.0 \\
\texttt{rtdl-revisiting-models} & 0.0.2 \\
\texttt{pytorch-tabnet} & 4.1.0 \\
\texttt{tabm} & 0.0.3 \\
\texttt{tabpfn} & 9.2.0 \\
\bottomrule
\end{tabularx}
\end{table}

Odtworzenie kampanii wymaga kolejno przygotowania środowiska z~zamrożonych zależności, audytu i~przygotowania danych, udostępnienia checkpointu oraz sprawdzenia GPU. Rewizję, nazwę pliku i~hash wag wskazuje konfiguracja. Powiązanie tych identyfikatorów z~manifestem wykonania rozróżnia modele korzystające z~tej samej wersji pakietu, lecz innych wag.

Manifesty, indeksy, historia prób i~predykcje tworzą podstawę odtworzenia raportów. Interpretacja naukowa pozostaje odrębna od działania raportera: wymaga określenia populacji porównania, znaczenia różnic i~zakresu niepewności.

\clearpage
\section{Wyniki badania porównawczego}
\label{chap:wyniki}

\textbf{Status rozdziału: szkielet do uzupełnienia po eksperymentach.} Nie zawiera wyników końcowej kampanii. Tabele są pustymi formularzami prezentacji, a~ramki wskazują docelowe ścieżki wykresów. Brak wartości nie oznacza zera, awarii ani niewykrycia różnicy. Po wykonaniu badania należy zastąpić instrukcje redakcyjne opisem obserwacji i~podać identyfikator uruchomienia oraz wersję analizy.

\subsection{Kompletność wykonania i odstępstwa}
\label{sec:wyniki-kompletnosc}

\textbf{Do uzupełnienia:} daty kampanii, wersje kodu i~danych oraz konfiguracja sprzętu oraz liczby wykonanych, udanych, nieudanych i~przerwanych prób. Punktem odniesienia są nominalne liczby z~sekcji~\ref{sec:koszt-metodyka}, a~nie założenie pełnego sukcesu. Tabelę~\ref{tab:wyniki-pokrycie} należy rozwinąć dla każdej pary zbiór--metoda. Osobno wskazać stały TabPFN, dla którego planowana jest jedna konfiguracja i~wyłącznie D.

\begin{table}[tbp]
\centering\small
\caption{Formularz kompletności kampanii; brak danych z~końcowego wykonania. Liczba kandydatów dotyczy różnych konfiguracji, a~nie liczby ponowionych prób.}
\label{tab:wyniki-pokrycie}
\begin{tabularx}{\linewidth}{YYYYY}
\toprule
Zbiór i~metoda & Planowane konfiguracje & Udane konfiguracje & Awarie i~przerwania & Dostępne warunki \\
\midrule
\multicolumn{5}{c}{Do uzupełnienia na podstawie indeksu, statusów i~historii prób.} \\
\bottomrule
\end{tabularx}
\end{table}

Należy opisać przyczyny braków, powtórzenia i~odstępstwa od zamrożonego protokołu. Nie zastępować braków arbitralnie niską metryką. Dla każdej analizy na podzbiorze kompletnych par podać jego skład oraz mianownik; wskazać, czy wykluczenia koncentrują się na trudniejszych zadaniach.

\subsection{Jakość na poszczególnych zadaniach}
\label{sec:wyniki-zadania}

\textbf{Do uzupełnienia:} wyniki wszystkich jedenastu procedur, bez wyboru przedstawiciela rodziny według testu. Panel główny obejmie T oraz odrębnie oznaczony stały TabPFN, zgodnie z~sekcją~\ref{sec:analiza-plan}. Tabele~\ref{tab:wyniki-binarne} i~\ref{tab:wyniki-wieloklasowe} celowo rozdzielają metryki o~różnych skalach i~kierunkach poprawy. Pełne D, T i~T+E oraz metryki pomocnicze należy udostępnić w~załączniku.

\begin{table}[tbp]
\centering\small
\caption{Formularz wyników dziewięciu zadań binarnych. ROC AUC: większa wartość oznacza lepszą jakość. Brak pomiarów.}
\label{tab:wyniki-binarne}
\begin{tabularx}{\linewidth}{YYYY}
\toprule
Zbiór & Metoda i~warunek & Testowe ROC AUC & Liczba obserwacji testowych \\
\midrule
\multicolumn{4}{c}{Do uzupełnienia z~wyrównanych predykcji i~raportu kampanii.} \\
\bottomrule
\end{tabularx}
\end{table}

\begin{table}[tbp]
\centering\small
\caption{Formularz wyników Covertype, Diabetes i~Heleny. Strata logarytmiczna: mniejsza wartość oznacza lepszą jakość. Brak pomiarów.}
\label{tab:wyniki-wieloklasowe}
\begin{tabularx}{\linewidth}{YYYY}
\toprule
Zbiór & Metoda i~warunek & Testowa strata logarytmiczna & Liczba obserwacji testowych \\
\midrule
\multicolumn{4}{c}{Do uzupełnienia; zachować mapowanie i~kolejność klas.} \\
\bottomrule
\end{tabularx}
\end{table}

\subsection{Uprzednio ustalone kontrasty}
\label{sec:wyniki-kontrasty}

\textbf{Do uzupełnienia po zamrożeniu analizy:} różnice RealMLP--CatBoost, TabM--CatBoost i~ModernNCA--KNN na każdym zadaniu, z~kierunkiem korzyści ustalonym w~sekcji~\ref{sec:analiza-plan}. Wskazać zakres niepewności, remisy, braki i~liczbę kompletnych porównań. Jeżeli zastosowano testy, podać wartości surowe i~skorygowane oraz uzasadnienie ich założeń. Nie interpretować braku istotności jako równoważności.

\plannedfigure
{figures/ch06/kontrasty.pdf}
{Miejsce na wykres uprzednio ustalonych kontrastów; wyniki nie zostały jeszcze uzyskane.}
{fig:wyniki-kontrasty}
{Oddzielne panele różnic AUC i~straty, trzy pary metod, dodatni kierunek korzyści pierwszej metody. Przedziały wyłącznie po zamrożeniu i~wykonaniu analizy.}

\subsection{Wpływ strojenia i łączenia konfiguracji}
\label{sec:wyniki-strojenie}

\textbf{Do uzupełnienia:} zmiana D--T i~T--T+E na teście, liczba skutecznych kandydatów oraz koszt ich znalezienia. Pokazać również pogorszenia. Oddzielić wynik wyboru na OOF od jego skutku testowego i~nie wybierać lepszego warunku po obejrzeniu testu. Nie tworzyć nieistniejących T i~T+E dla TabPFN.

\plannedfigure
{figures/ch06/strojenie-i-zespoly.pdf}
{Miejsce na porównanie D, T i~T+E; wyniki nie zostały jeszcze uzyskane.}
{fig:wyniki-strojenie}
{Pokazać sparowane zmiany na tych samych zadaniach. Rozdzielić AUC i~stratę, wskazać liczbę udanych konfiguracji i~jawnie oznaczyć wyjątek TabPFN.}

\subsection{Zróżnicowanie wyników między zadaniami}
\label{sec:wyniki-heterogenicznosc}

\textbf{Do uzupełnienia:} opis wyników w~odniesieniu do audytowanych właściwości danych. Przypadki małe, szerokie, kategoryczne, niezrównoważone i~wieloklasowe mogą się nakładać; nie są niezależnymi grupami kontrolnymi. Oddzielić obserwowany wzorzec od hipotezy o~jego przyczynie. Dwanaście punktów nie uzasadnia rozbudowanego modelu wyjaśniającego przewagę architektur. Uwzględnić wyniki sprzeczne z~oczekiwaniami.

\subsection{Koszt i wykonalność}
\label{sec:wyniki-koszt}

\textbf{Do uzupełnienia:} pomiary zgodne z~granicami sekcji~\ref{sec:koszt-metodyka}. W~tabeli~\ref{tab:wyniki-koszt} czas wyszukiwania, wybranych konfiguracji i~predykcji pozostają odrębne. Zarejestrowaną sumę przebiegów predykcji należy tak nazwać; pełną latencję zespołu można podać tylko po jej niezależnym pomiarze. Pamięć i~awarie zasobów wymagają identyfikowalnego źródła, a~brak odczytu nie oznacza braku zużycia. Obecne pomiary pamięci pochodzą z~rekordów konfiguracji i~obejmują dopasowanie, nie testową predykcję. Nie należy przedstawiać ich jako szczytu całej próby; osobne pomiary inferencji wymagają wskazania ich artefaktów.

\begin{table}[tbp]
\centering\small
\caption{Formularz kosztów procedur; brak pomiarów. Jednostką czasu będą sekundy, a~pamięci --- jawnie wskazane MiB lub GiB.}
\label{tab:wyniki-koszt}
\begin{tabularx}{\linewidth}{YYYYY}
\toprule
Zadanie, metoda i~warunek & Koszt wyszukiwania & Dopasowanie wybranych konfiguracji & Zapisany koszt predykcji & Sprzęt i~pamięć \\
\midrule
\multicolumn{5}{c}{Do uzupełnienia z~artefaktów i~zweryfikowanych pomiarów zasobów.} \\
\bottomrule
\end{tabularx}
\end{table}

\plannedfigure
{figures/ch06/jakosc-koszt.pdf}
{Miejsce na wykres kompromisu jakości i~kosztu w~obrębie zadań; brak pomiarów.}
{fig:wyniki-koszt}
{Osobny panel dla zadania. Zdefiniować oś kosztu i~warunek, oznaczyć CPU/GPU. Nie sumować kosztów D, T i~T+E jako niezależnych wyszukiwań.}

\subsection{Odpowiedzi na pytania PB1, PB2 i PB5}
\label{sec:wyniki-odpowiedzi}

\textbf{Do uzupełnienia po analizie:} trzy osobne odpowiedzi oparte na numerach tabel i~rysunków, wielkości efektów oraz pokryciu. Dla PB1 określić, w~jakich zadaniach i~warunkach wystąpiły różnice; dla PB2 oddzielić strojenie od łączenia konfiguracji; dla PB5 wskazać jakość wraz z~kosztem i~ograniczeniami sprzętowymi. Interpretacja literaturowa i~ogólne ograniczenia należą do końcowej dyskusji.

\chapter{Badania mechanizmów uwrażliwiania na rzadkie klasy i ograniczania nadmiernego dopasowania}
\label{chap:mechanizmy}

\textbf{Status rozdziału: szkielet do uzupełnienia po implementacji i~wykonaniu badań A oraz B.} Nie zawiera wyników interwencji ani dowodu ich skuteczności. Projekt opisano w~sekcjach~\ref{sec:projekt-wazenie} i~\ref{sec:projekt-regularyzacja}. Poniższe instrukcje określają wymagany materiał wynikowy; nie zastępują jeszcze niezamrożonych ustawień.

\section{Wykonane warunki i zgodność kontroli}
\label{sec:mechanizmy-kontrole}

\textbf{Do uzupełnienia:} identyfikatory badań, kompletność warunków, trzy ziarna uczenia, rzeczywiste indeksy foldów i~podzbiorów, etykiety rzadkie oraz wybrane konfiguracje. Dla A potwierdzić wyliczenie wag wyłącznie z~treningu i~brak innych zmian hiperparametrów. Dla B podać pochodzenie niezależnej receptury, wartości dropoutu i~zaniku wag, liczbę przykładów każdej klasy, reguły zatrzymania i~wyboru epoki. Każde odstępstwo trzeba opisać przed interpretacją efektu. Wynik bez dowodu dopasowania kontroli nie jest izolowaną oceną interwencji.

\section{Ważenie klas a jakość uporządkowania}
\label{sec:mechanizmy-ap}

\textbf{Do uzupełnienia:} AP rzadkiej klasy dla obu warunków, osobno dla zbioru, metody i~ziarna. Tabelę~\ref{tab:mechanizmy-ap} należy uzupełnić wynikami sparowanymi, a~następnie średnią różnicą i~zakresem trzech różnic. Podać surową etykietę rzadką, częstość i~liczbę jej przykładów testowych. ROC AUC oraz strata logarytmiczna pozostają ocenami uzupełniającymi; poprawa jednej miary nie przesądza o~pozostałych.

\begin{table}[tbp]
\centering\small
\caption{Formularz efektu ważenia na AP klasy rzadkiej. Brak wyników; dodatnia różnica ma oznaczać korzyść ważenia.}
\label{tab:mechanizmy-ap}
\begin{tabularx}{\linewidth}{YYYYY}
\toprule
Zbiór i~metoda & Ziarno & AP bez wag & AP z~wagami & Sparowana różnica \\
\midrule
\multicolumn{5}{c}{Do uzupełnienia z~predykcji testowych badania A.} \\
\bottomrule
\end{tabularx}
\end{table}

\plannedfigure
{figures/ch07/precision-recall.pdf}
{Miejsce na testowe krzywe precision--recall badania ważenia; brak wyników.}
{fig:mechanizmy-pr}
{Pokazać obie interwencje i~wszystkie trzy ziarna. Zdefiniować klasę rzadką; zaznaczyć jej częstość. AP liczyć z~predykcji, bez utożsamiania z~trapezowym polem.}

\section{Ważenie, próg i punkt pracy}
\label{sec:mechanizmy-progi}

\textbf{Do uzupełnienia:} zamrożona reguła progowania, progi wybrane na OOF oraz testowe precyzja, czułość, FP, FN i~udział alarmów. Należy potwierdzić, że etykiety testowe nie uczestniczyły w~wyborze. W~tabeli~\ref{tab:mechanizmy-progi} każda komórka musi identyfikować warunek i~ziarno. Jeżeli punkt pracy poprawił czułość kosztem liczby alarmów, należy przedstawić obie zmiany. Zmiana progu i~zmiana jakości uporządkowania są odrębnymi wynikami.

\begin{table}[tbp]
\centering\small
\caption{Formularz oceny progów wybranych na OOF; brak wyników. Miary decyzyjne dotyczą wyłącznie zewnętrznego testu.}
\label{tab:mechanizmy-progi}
\begin{tabularx}{\linewidth}{P{0.20\linewidth}YYYYY}
\toprule
Zadanie, metoda, warunek i~ziarno & Próg OOF & Precyzja & Czułość & FP i~FN & Udział alarmów \\
\midrule
\multicolumn{6}{c}{Do uzupełnienia po zapisaniu wyboru progów i~ocenie testowej.} \\
\bottomrule
\end{tabularx}
\end{table}

W~interpretacji trzeba uwzględnić liczebność rzadkich zdarzeń, ewentualną zmianę jakości prawdopodobieństw i~ograniczenia przenoszenia progu z~OOF na średnią modeli foldów. Bez zadanych kosztów błędów nie można uznać konkretnego punktu pracy za optymalny wdrożeniowo.

\section{Liczność treningu i ablacje regularyzacji MLP}
\label{sec:mechanizmy-licznosc}

\textbf{Do uzupełnienia:} wyniki pełnej receptury, wariantu bez dropoutu i~wariantu bez zaniku wag przy 25\%, 50\% i~100\% treningu. W~tabeli~\ref{tab:mechanizmy-ablacje} należy pokazać każde ziarno, rzeczywiste liczności treningowe i~różnice względem pełnej receptury dla tych samych indeksów. Potwierdzić zagnieżdżenie podzbiorów oraz niezmienność walidacji i~testu. Liczności poszczególnych foldów można rozwinąć w~załączniku; nie wolno zastąpić ich jedną pozornie wspólną liczbą.

\begin{table}[tbp]
\centering\small
\caption{Formularz wyników badania B. Brak wyników; nominalny udział odnosi się do treningu wewnętrznego foldu.}
\label{tab:mechanizmy-ablacje}
\begin{tabularx}{\linewidth}{YYYYY}
\toprule
Zbiór, udział i~ziarno & Rzeczywiste liczności & Wariant receptury & Testowe AUC & Różnica względem pełnej receptury \\
\midrule
\multicolumn{5}{c}{Do uzupełnienia z~manifestów podzbiorów i~predykcji badania B.} \\
\bottomrule
\end{tabularx}
\end{table}

\plannedfigure
{figures/ch07/jakosc-a-licznosc.pdf}
{Miejsce na zależność jakości od liczby przykładów i~wariantu regularyzacji; brak wyników.}
{fig:mechanizmy-licznosc}
{Osobne panele German Credit, Ionosphere i~Bioresponse; trzy receptury i~trzy ziarna. Nie przedstawiać zakresu trzech ziaren jako przedziału ufności.}

\section{Dynamika uczenia i luka generalizacji}
\label{sec:mechanizmy-luki}

\textbf{Do uzupełnienia:} historie nieważonej straty treningowej i~walidacyjnej tej samej sieci, obliczone w~trybie ewaluacji. Podać definicję luki, wybrane epoki i~regułę agregowania. Krzywe należy oprzeć na zapisanych epokach, a~nie na interpolowanych wynikach końcowych. Nie wybierać do prezentacji wyłącznie ziarna z~największą różnicą; selekcja paneli musi zostać uzasadniona i~zapewniać dostęp do pozostałych przebiegów.

\plannedfigure
{figures/ch07/krzywe-i-luki.pdf}
{Miejsce na krzywe uczenia i~luki trening--walidacja badania B; brak wyników.}
{fig:mechanizmy-luki}
{Porównywać tę samą nieważoną stratę dla jednej sieci foldu z~wyłączonym dropoutem. Podać fold, ziarno, podzbiór i~epokę. Nie łączyć treningu sieci z~testem zespołu.}

Zmniejszenie luki należy ocenić razem z~poziomem straty i~jakością na teście. Pogorszenie dopasowania treningowego przy niezmienionym lub gorszym teście nie jest dowodem skutecznego ograniczenia nadmiernego dopasowania w~znaczeniu poprawy generalizacji. Zróżnicowanie między ziarnami musi pozostać widoczne.

\section{Odpowiedzi na pytania PB3 i PB4}
\label{sec:mechanizmy-odpowiedzi}

\textbf{Do uzupełnienia po analizie:} osobne odpowiedzi określające kierunek, wielkość i~warunki obserwowanego efektu. PB3 wymaga rozdzielenia AP, punktu pracy i~jakości prawdopodobieństw. PB4 wymaga powiązania ablacji z~licznością treningu, zmianą jakości i~luką. Należy uwzględnić przypadki braku poprawy i~przeciwne kierunki między zadaniami. Zakres odpowiedzi obejmuje po trzy celowo wybrane zbiory, ustalone receptury i~trzy ziarna; nie jest ogólnym twierdzeniem o~ważeniu lub regularyzacji dla wszystkich tabel.


\printbibliography[heading=bibintoc,title={Bibliografia}]
\end{document}
